\documentclass[11pt,reqno]{amsart}
\usepackage[T1]{fontenc}
\usepackage[utf8]{inputenc}
\usepackage{lmodern}
\usepackage{amsmath,amssymb,amsthm,mathtools}
\usepackage{booktabs,longtable,array}
\usepackage{xurl}
\usepackage[hidelinks,hypertexnames=false]{hyperref}
\usepackage{microtype}
\allowdisplaybreaks
\setlength{\emergencystretch}{2em}
\newtheorem{theorem}{Theorem}[section]
\newtheorem{lemma}[theorem]{Lemma}
\newtheorem{proposition}[theorem]{Proposition}
\newtheorem{corollary}[theorem]{Corollary}
\theoremstyle{definition}
\newtheorem{definition}[theorem]{Definition}
\theoremstyle{remark}
\newtheorem{remark}[theorem]{Remark}
\newcommand{\R}{\mathbb R}
\newcommand{\Q}{\mathbb Q}
\newcommand{\mx}[1]{\left\|#1\right\|_{\max}}
\newcommand{\iin}[1]{\left\|#1\right\|_\infty}
\newcommand{\diag}{\operatorname{diag}}
\newcommand{\amin}{\alpha_-}
\newcommand{\gam}{\gamma}
\newcommand{\cB}{\mathcal B}
\DeclareMathOperator{\sgn}{sgn}
\title[The exact growth factor for complete pivoting]
{The exact maximum growth factor for\\
complete pivoting on real \(5\times5\) matrices}
% Author order updated by the user on 12 September 2026: 马千里, 吴超.
\author{Qianli Ma}
\address[Qianli Ma]{Zhejiang University; WUJIE AI}
\email{qianli.ma@zju.edu.cn}
\author{Chao Wu}
\address[Chao Wu]{Zhejiang University}
\email{chao.wu@zju.edu.cn}
\hypersetup{pdfauthor={Qianli Ma, Chao Wu},
pdftitle={The exact maximum growth factor for complete pivoting on real 5 x 5 matrices}}
\date{11 September 2026}
\begin{document}

\begin{abstract}
Let \(\rho_5^{\R}\) be the largest element growth possible in exact
Gaussian elimination with complete pivoting on real \(5\times5\) matrices,
with every legal choice at a pivot tie allowed. We prove
\[
\rho_5^{\R}=\alpha
=4.132517078632472854223346853277\ldots,
\]
where \(\alpha\) is the algebraic lower bound identified by
Chen, Edelman, and Urschel, specified by their degree-\(61\) polynomial and
a rational isolating interval. Our contribution is the global upper bound.
An exact inverse-elimination representation retains all initial and
intermediate entry bounds of the same matrix. A constructive contraction
preserves the final-pivot height while reducing the problem to two
saturated tail geometries. On the negative-diagonal geometry, three
feasible paths eliminate six variables and give an attained fibre maximum
over seventeen shared parameters. Rational certificates control the
remaining domains, while a constrained algebraic maximum argument and
feasible transports resolve regions meeting the sharp value. The proof
combines these analytic reductions with exact verification of a finite,
complete cover. We provide a partial Lean formalization together with publicly
available exact certificates and verification code.
\end{abstract}

\maketitle

\section{The problem and the result}

Complete pivoting selects a largest-modulus entry of the current Schur
complement at every step of Gaussian elimination. Although this rule
controls local multipliers, successive Schur updates can generate entries
larger than those in the input. Sharp growth factors in small dimensions
provide concrete tests of what global information is retained, or lost,
by general pivot-growth estimates.

The values in dimensions three and four are \(9/4\) and \(4\),
respectively; classical work includes Cryer~\cite{Cryer1968} and
Cohen~\cite{Cohen1974}. Numerical evidence for a fifth-pivot value near
\(4.1325\) was reported by Day and Peterson~\cite{DayPeterson1988}
and subsequently reproduced by Gould~\cite{Gould1991}.
Chen, Edelman, and Urschel~\cite{ChenEdelmanUrschel2026} identified the
specific degree-\(61\) algebraic value \(\alpha\) and proved a lower bound
by a completely pivoted matrix. Their Theorem~2.2 proves optimality within
a prescribed equality family and supplies the lower-bound direction used
here; their Theorem~3.5 gives the global upper bound \(4.84\).
Theorem~\ref{thm:main} establishes the equality in their Conjecture~2.1.
The constant, its defining polynomial, and the equality pattern are
due to Chen, Edelman, and Urschel; the advance here is the matching
global upper bound.

The present result concerns a fixed dimension and the real field.
It is separate from estimates as the dimension tends to infinity,
including the recent results of
Bisain, Edelman, and Urschel~\cite{BisainEdelmanUrschel2025} and
Shah and Urschel~\cite{ShahUrschel2026}.
Background on growth and its relation to numerical error can be found in
Day and Peterson~\cite{DayPeterson1988} and
Higham and Higham~\cite{HighamHigham1989}.

\begin{definition}
For \(0\ne A\in\R^{n\times n}\), let \(\pi\) denote a legal complete-pivoting
path in exact arithmetic. Write \(A^{(j)}_\pi\) for the trailing Schur
matrix at stage \(j\), before its leading pivot is eliminated. Define
\[
\rho(A,\pi)=
\frac{\max_j\mx{A^{(j)}_\pi}}{\mx A},
\qquad
\rho_n^\R=\sup_{A\ne0}\ \sup_{\pi\ {\rm legal}}\rho(A,\pi).
\]
If the current pivot is zero, the trailing matrix is zero and elimination
terminates; all subsequent pivot magnitudes are taken to be zero.
\end{definition}

We normalize \(\mx A=1\). The first pivot magnitude is then one, and
complete pivoting gives
\begin{equation}\label{eq:growthpivots}
\rho(A,\pi)=\max_{1\le j\le n}p_j,
\end{equation}
where \(p_j\) is the magnitude of the chosen pivot. Row and column
permutations may be absorbed into the original matrix. Thus it suffices
to prove the upper bound for every matrix that is completely pivoted in
the fixed diagonal order. This does not assume uniqueness of a largest
entry, nor equality of the results of different tied paths.

\begin{theorem}\label{thm:main}
For exact complete-pivoting Gaussian elimination on real \(5\times5\)
matrices, \(\rho_5^\R=\alpha\).
Here \(\alpha\) is the root specified in
Section~\ref{sec:alpha}. The upper bound holds for every legal pivot path,
including paths with ties and inputs for which elimination terminates
early.
\end{theorem}

% Included section: readability_main
% New mathematical reading guide; integrated after the main theorem.
\subsection*{Why the algebraic candidate does not settle the problem}
An attaining matrix supplies one legal elimination path with growth
\(\alpha\). Even a proof that this matrix is optimal among all matrices
with the same active equalities leaves other equality patterns and
interior feasible points to be controlled. This is the distinction between
the restricted-family result of
\cite[Theorem~2.2]{ChenEdelmanUrschel2026} and the global conclusion here.
Their work also gives a global upper bound of \(4.84\), using interval
covers; the question addressed in the present paper is how to reach the
sharp value without assuming the candidate's equality pattern.

A second obstacle is that the successive Schur complements are coupled.
An upper bound on one entry may be attained by a different completion
from the one attaining another bound. Combining such separate extrema
can create a tail that no input matrix produces. We retain the entire
input matrix in an inverse-elimination representation, and every analytic
motion below is a motion of that same complete matrix. In particular,
the six-variable maximization is attained by a feasible completion;
it is not a collection of independently optimized capacities.

The sharp value presents a further difficulty for strict numerical
exclusion. Let \(\gamma<\alpha\) be the rational threshold fixed below.
A region containing an attaining matrix cannot be ruled out under the
hypothesis \(F\ge\gamma\). Nor can a strict rational estimate
\(F\le\alpha_-<\alpha\) cover that matrix. Such regions require a
separate exact-height argument or a feasible transport to a domain
already bounded by \(\alpha\). This explains the two types of terminal
used in the proof: some exclude high points, while others retain feasible
points and prove that they are safe at the sharp value.

\subsection*{The architecture of the upper bound}
The proof is organized around the following reductions. Here \(F\) is
the magnitude of the fifth pivot of a normalized matrix.
\begin{enumerate}
\item \emph{Remove earlier-stage growth.}
The bounds \(1,2,9/4,4\) on the first four pivots hold along every
legal path, including ties and early termination. A counterexample to
the theorem must therefore have \(F>\alpha\).
\item \emph{Reach a saturated tail without losing height.}
Two contractions of actual rows and columns send every positive
fifth-pivot value to a diagonal-saturated or cross-saturated tail.
Signs and tied boundary cases reduce these to the complete \(X\) domain
or the proper negative-diagonal domain. No rank or generic-position
assumption is introduced into this final division of the problem.
\item \emph{Control the cross-saturated domain.}
Four overlapping parameter ranges cover all of \(X\). The ranges away
from the peak admit rational exclusions. Near the peak, a canonical
representative theorem, finite localization, and a constrained algebraic
maximum argument give the exact bound. This result is established before
it is used in the negative-diagonal proof.
\item \emph{Maximize a complete negative-diagonal fibre.}
At fixed seventeen shared parameters, three feasible paths maximize six
remaining variables. Their endpoint realizes the maximum of the whole
fibre. The question is thereby reduced to bounding one explicitly defined
capacity on the space of complete frameworks.
\item \emph{Cover the remaining parameter domain.}
Finite certificates use exact subdivisions, necessary inequalities,
qualified contractions, and the two terminal types described above.
A complete source is followed through these operations to a justified
terminal. The original-root composition accounts for every residual
branch, including branches inherited through changes of coordinates.
\end{enumerate}

The finite computations carry specific parts of Steps 3 and 5. The
written argument explains both the domains they must cover and why their
acceptance rules imply the required real inequalities. The supplement
specifies the exact proof objects used for those obligations. One can
therefore audit the reduction and the verifier logic from the text;
checking every large proof object additionally requires the supplied
exact-arithmetic programs or an independent implementation of their rules.

\begin{center}
\small
\begin{tabular}{@{}>{\raggedright\arraybackslash}p{.25\linewidth}>{\raggedright\arraybackslash}p{.32\linewidth}>{\raggedright\arraybackslash}p{.33\linewidth}@{}}
\toprule
Mathematical obligation&Reduction established in the text&Finite evidence consumed\\
\midrule
Arbitrary legal real input&Early-pivot bounds, complete lift, same-height saturation&None for these reductions\\
Complete \(X\) domain&Four-range union; qualified canonical and local height rules&Exclusion and localization certificates on the stated sources\\
Negative-diagonal fibre&An attained six-variable maximum with full physical bands&None for the fibre formula\\
All high frameworks&Source-preserving rules and mixed-height coverage&Complete root, continuation, and safety certificates\\
Sharp value&Upper bound together with the known attained lower bound&Exact identities binding the algebraic constant\\
\bottomrule
\end{tabular}
\end{center}


\section{The constant and an attained lower bound}\label{sec:alpha}

Let
\begin{equation}\label{eq:P61}
P_{61}(g)=\sum_{i=0}^{61}c_i g^i
\end{equation}
be the polynomial in Chen, Edelman, and
Urschel~\cite[Eq.~(2.15)]{ChenEdelmanUrschel2026}. Its leading coefficient is \(59049\)
and its constant coefficient is
\[
-400768351683901408958416200638054732473753927680.
\]
Supplementary Material~S1 gives the full ascending coefficient list
and exact decimal rational endpoints \(L<U\), with
\(U-L=10^{-200}\), isolating one real root. Throughout, \(\alpha\) means this root and \(\amin=L\).
There is also a short exact identification with the published constant.
All \(62\) integer coefficients agree with the published polynomial.
For
\[
Q(t)=(1+t)^{61}P_{61}\!\left(\frac{4+5t}{1+t}\right),
\]
the coefficients of degrees \(0,\ldots,8\) are negative and those
of degrees \(9,\ldots,61\) are positive.
Descartes' rule of signs therefore gives at most one positive root
of \(Q\), hence at most one root of \(P_{61}\) in \((4,5)\).
Exact evaluation gives \(P_{61}(L)<0<P_{61}(U)\).
Consequently \(\alpha\) is the unique root in \((4,5)\).
The exact coefficient comparison, transformed coefficients, and
endpoint signs are recorded in Supplementary Material~S1.
For orientation,
\[
\begin{gathered}
4.132517078632472854223346853277<\alpha,\\
\alpha<4.132517078632472854223346853278.
\end{gathered}
\]
All certificate comparisons use the full rational endpoints, not these
truncated decimals. We also use
\begin{equation}\label{eq:thresholds}
\Lambda=\frac{1653}{400},\qquad
\gam=\frac{4132517}{10^6},\qquad
4<\Lambda<\gam<\amin<\alpha.
\end{equation}

The lower bound and an attaining completely pivoted matrix are
\cite[Theorem~2.2]{ChenEdelmanUrschel2026}. The coefficient comparison
and root identification above identify their constant with the present
\(\alpha\). Only this established lower-bound direction is needed here.

Supplementary Material~S1 also identifies an independent exact
reconstruction of the same algebraic candidate. Its archived equations,
rational reconstruction functions, nonzero-denominator checks, and
complete-pivoting inequalities allow the witness to be checked separately.
We keep those longer calculations with the supplementary data so that
the main text can develop the new global upper bound. Neither the
candidate's existence nor the root isolation alone gives that upper bound.

\section{An elementary reduction to the fifth pivot}\label{sec:scalar}

To bound the growth of a five-dimensional elimination, it is not enough
to bound its fifth pivot unless all earlier pivots are already safe.
We therefore give the elementary bounds used both for that reduction
and for the parameter domains later in the proof. They also settle
paths that terminate before the fifth stage.

\begin{lemma}\label{lem:Phi}
For a normalized real \(3\times3\) matrix completely pivoted in
the indicated order, let its pivot magnitudes be \(1,t,g\). Then
\[
0\le t\le2,\qquad g\le\Phi(t),\qquad
\Phi(t)=
\begin{cases}
2t,&0\le t\le1,\\
t(3-t),&1\le t\le2.
\end{cases}
\]
Each branch is sharp. In particular, \(g\le9/4\).
\end{lemma}

\begin{proof}
The first Schur update is a difference of two quantities of magnitude at
most one, so \(t\le2\). If \(t=0\), its whole trailing block is zero.
Otherwise the last update gives \(g\le2t\), proving the first branch.

For \(1<t\le2\), normalize the first row and column to be nonnegative by
independent sign changes. A positive entry of the first Schur complement
is at most one, so the second pivot is negative. Write the tail as
\[
\begin{pmatrix}-t&\gamma_0\\\eta_0&\zeta_0\end{pmatrix}.
\]
Let the relevant first-row and first-column coordinates be
\(x,u,y,v\in[0,1]\), and set
\[
c=t-1,\quad s=xv,\quad z=uy,\quad r=uv,\quad w=xy.
\]
The original entry at the pivot is \(w-t\), whence \(w\ge c\).
Furthermore \(sz=wr\) and \(r\ge sz\). Intersecting the original entry
bounds with the Schur pivot bounds gives
\begin{equation}\label{eq:threeintervals}
\begin{split}
\gamma_0&\in[-1-\min(c,s),\,1-s],\\
\eta_0&\in[-1-\min(c,z),\,1-z],\\
\zeta_0&\in[-1-\min(c,r),\,1-r].
\end{split}
\end{equation}
We claim that, with \(M=t^2(3-t)\),
\begin{align}
t(1-r)+(1+\min(c,s))(1+\min(c,z))&\le M,\label{eq:aux1}\\
t(1+\min(c,r))+(1+\min(c,s))(1-z)&\le M.\label{eq:aux2}
\end{align}
For \eqref{eq:aux1}, replace \(r\) by \(sz\). When \(s,z\le c\)
the resulting function increases in either variable, since its
corresponding derivative is \(1-cz\) or \(1-cs\); its maximum is at
\(s=z=c\). If \(s\le c\le z\), it is
\(t[2+s(1-z)]\le t[2+c(1-c)]=M\).
The other mixed case is symmetric. If \(s,z\ge c\), the bound is
\(t(1-c^2)+t^2=M\).

For \eqref{eq:aux2}, the shared-product relations imply
\(z\ge cr\), and also \(z\ge r\) whenever \(s\le c\).
The latter includes \(s=0\): then \(sz=wr\) and \(w\ge c>0\)
force \(r=0\). Thus no division by \(s\) is needed.
Separate four closed cases.
If \(r\ge c,s\ge c\), then \(z\ge cr\ge c^2\), and the expression is at
most \(t^2+t(1-c^2)=M\).
If \(r\ge c,s\le c\), then \(z\ge r\ge c\), and the bound
\(t^2+t(1-c)=2t\le M\) suffices.
If \(r\le c,s\le c\), the expression is at most
\[
t(1+r)+(1+s)(1-r)\le2t\le M;
\]
the first difference is \((c-s)(1-r)\ge0\), and
\(M-2t=tc(1-c)\ge0\).
Finally, if \(r\le c,s\ge c\), then
\(r-z\le r(1-c)\le c(1-c)\).
The expression is at most \(t[2+c(1-c)]=M\).
Interchanging \(s,z\) gives the symmetric form of \eqref{eq:aux2}.

Put \(A_s=1+\min(c,s)\), \(B_s=1-s\), and similarly for \(z,r\).
By \eqref{eq:threeintervals},
\[
t\zeta_0+\gamma_0\eta_0
\le tB_r+\max(A_sA_z,B_sB_z)\le M.
\]
The first product uses \eqref{eq:aux1}; the second is at most
\(t+1\le M\). On the other side,
\[
-(t\zeta_0+\gamma_0\eta_0)
\le tA_r+\max(A_sB_z,B_sA_z)\le M
\]
by \eqref{eq:aux2} and its symmetric version.
Thus the magnitude of the tail determinant is at most \(M\), and
division by \(t\) proves \(g\le t(3-t)\).

Sharp examples are
\[
\begin{pmatrix}1&0&0\\0&t&t\\0&-t&t\end{pmatrix}
\quad(0\le t\le1),\qquad
\begin{pmatrix}1&1&c\\1&-c&-1\\c&-1&1\end{pmatrix},
\quad c=t-1\quad(1\le t\le2).
\]
Finally \(9/4-t(3-t)=(t-3/2)^2\).
\end{proof}

\begin{corollary}\label{cor:firstfour}
The first four pivot magnitudes of a normalized real completely pivoted
matrix of order at most five satisfy
\[
p_1=1,\qquad p_2\le2,\qquad p_3\le9/4,\qquad p_4\le4.
\]
Consequently \(\rho(A,\pi)>\alpha\) implies \(p_5>\alpha\).
\end{corollary}

\begin{proof}
Any leading submatrix in the chosen pivot order is completely
pivoted, and a Schur tail divided by its first pivot magnitude is
normalized and completely pivoted. Put \(x=p_2\), \(y=p_3/p_2\)
when these denominators are nonzero. Lemma~\ref{lem:Phi} gives
\[
xy\le\Phi(x),\qquad p_4\le x\Phi(y).
\]
If \(x\le1\), the second bound is at most \(9/4\). If
\(1\le x\le3/2\), it is at most \(27/8\). If \(3/2\le x\le2\),
then \(y\le3-x\le3/2\). The function \(\Phi\) is nondecreasing on
\([0,3/2]\), so
\[
p_4\le x\Phi(3-x)=x^2(3-x)\le4,
\]
using \(4-x^2(3-x)=(2-x)^2(x+1)\).
Zero pivots imply zero subsequent growth.
Equation~\eqref{eq:growthpivots} and \(\alpha>4\) finish the proof.
\end{proof}

\section{A complete two-step lift}\label{sec:lift}

The representation in this section is the point at which entry bounds
from different stages are coupled. Starting from the third Schur block,
we undo the preceding two elimination steps and retain the original
matrix as well. Consequently every row, column, and auxiliary product
used later has a simultaneous physical interpretation.

All subsequent variables describe one common matrix. Let
\(u,x,v,q\in\R^3\), \(p,k>0\), and \(e,\beta\in\R\).
Set
\begin{equation}\label{eq:lift}
M=
\begin{pmatrix}
1&-e&v^T\\
\beta&p-e\beta&(q+\beta v)^T\\
u&px-eu&D+xq^T+uv^T
\end{pmatrix},\qquad D=kC,\quad C_{00}=1.
\end{equation}
Subscripts on \(3\times3\) blocks run from zero to two.
Its first two Schur tails are
\begin{equation}\label{eq:schurtails}
M^{(2)}=
\begin{pmatrix}p&q^T\\px&D+xq^T\end{pmatrix},
\qquad M^{(3)}=D.
\end{equation}

\begin{proposition}\label{prop:lift}
Suppose \(C\) is completely pivoted in the fixed order, with
\(\mx C=1\). Then \eqref{eq:lift} is normalized and completely pivoted,
with first three pivot magnitudes \(1,p,k\), if and only if
\begin{align}
|e|,|\beta|,\iin u,\iin v,\iin x&\le1,\label{eq:bands1}\\
|p-e\beta|,\ \iin{px-eu},\ \iin{q+\beta v}&\le1,\label{eq:bands2}\\
\iin q\le p,\qquad \mx{D+xq^T}&\le p,\label{eq:bands3}\\
\mx{D+xq^T+uv^T}&\le1.\label{eq:bands4}
\end{align}
Every normalized fixed-order completely pivoted matrix with three
nonzero leading pivots has this representation after row-sign
normalization of those pivots.
\end{proposition}

\begin{proof}
The identities \eqref{eq:schurtails} show sufficiency: the first matrix
has entries bounded by one, the first tail by \(p\), and the second tail
is \(kC\). Conversely recover \(e,\beta,u,v\) from the initial matrix,
then \(p,x,q\) from its first Schur tail. The remaining tail is uniquely
\(D\); division by its positive leading pivot gives \(C\).
The inequalities are precisely the initial and first-stage bounds.
No independent completion is chosen for different inequalities.
\end{proof}

The alternative notation \(q=py\) makes the bottom block
\(kC+uv^T+pxy^T\). The actual variable \(q\), rather than \(y\), will be
used in the seventeen-dimensional reduction.
Write
\begin{equation}\label{eq:core}
C=
\begin{pmatrix}
1&a&b\\
c&ca+\rho&cb+\sigma\\
d&da+\tau&db+\omega
\end{pmatrix}.
\end{equation}
The fixed-order core conditions are
\(\mx C\le1\) and \(|\sigma|,|\tau|,|\omega|\le|\rho|\).
For \(\rho\ne0\), set
\[
h(C)=\omega-\tau\sigma/\rho,\qquad F=k|h(C)|.
\]
If \(\rho=0\), the final tail is zero; set \(h(C)=F=0\).
Lemma~\ref{lem:Phi} also yields
\begin{equation}\label{eq:necessaryscalars}
k\le9/4,\quad |h(C)|\le9/4,\quad k\le2p,\quad
|h(C)|\le2|\rho|.
\end{equation}
All these quantities belong to the same source.

\section{Constructive saturation of the last tail}\label{sec:saturation}

The next reduction explains why two tail geometries suffice. It constructs
a new complete matrix with the same last-pivot height from each input;
it does not infer the candidate's tight constraints from numerical
experiments. The construction uses contractions of actual rows and
columns, which lets us check all earlier pivoting conditions at once.

The next reduction is an existence statement about representatives.
It does not assert that every maximizer is already saturated.

\begin{lemma}[Same-height diagonal saturation]\label{lem:saturation}
Every complete lift \eqref{eq:lift} with \(F>0\) has another complete
lift with the same \(p,k,F\) whose normalized tail satisfies either
\[
D:\quad |\omega|=|\rho|,\quad |\sigma|,|\tau|\le|\rho|,
\]
or
\[
X:\quad |\sigma|=|\tau|=|\rho|,\quad |\omega|\le|\rho|.
\]
The construction is valid on every legal tied pivot path.
\end{lemma}

\begin{proof}
For \(0\le\lambda,\mu\le1\), let
\(R_\lambda=\diag(1,\lambda,1)\), \(K_\mu=\diag(1,\mu,1)\).
Replace
\[
C\mapsto R_\lambda C K_\mu,\quad
(u,x)\mapsto(R_\lambda u,R_\lambda x),\quad
(v,q)\mapsto(K_\mu v,K_\mu q).
\]
The new full matrix is
\(\diag(I_2,R_\lambda)M\diag(I_2,K_\mu)\).
Its initial and first-tail entries are only contracted in magnitude,
while the pivots \(1,p,k\) do not change.
The core tail transforms as
\begin{equation}\label{eq:sattransform}
(\rho,\sigma,\tau,\omega)
\mapsto(\lambda\mu\rho,\lambda\sigma,\mu\tau,\omega).
\end{equation}
Since \(F>0\), \(\rho\ne0\). Put
\[
S_0=|\sigma/\rho|,\quad T_0=|\tau/\rho|,\quad
W_0=|\omega/\rho|,\qquad S_0,T_0,W_0\in[0,1].
\]
If \(W_0\ge S_0T_0\), choose
\(\lambda=\max(T_0,W_0)\), \(\mu=W_0/\lambda\).
If \(W_0\) vanished, then \(S_0T_0=0\) and the original height would
be zero. Thus the division is valid and \(\lambda\mu>0\).
Now \(\lambda\ge T_0\), \(\mu\ge S_0\), and
\(\lambda\mu=W_0\), so the new completely pivoted tail lies on \(D\).
If \(W_0\le S_0T_0\), take \(\lambda=T_0,\mu=S_0\).
Positivity of the original height again excludes \(\lambda\mu=0\);
the new tail lies on \(X\). At the positive-product junction
\(W_0=S_0T_0>0\), the constructions agree.
Finally,
\[
\omega-\frac{(\mu\tau)(\lambda\sigma)}{\lambda\mu\rho}
=\omega-\frac{\tau\sigma}{\rho}.
\]
This preserves \(F\) and all the complete early-stage conditions.
No exchange of unequal pivots is introduced.
\end{proof}

Zero fifth pivots are already safe by Corollary~\ref{cor:firstfour}.
They need not have a zero fourth-stage tail: a nonzero rank-one
\(2\times2\) tail can have zero last pivot. In contrast, a zero
fourth pivot under complete pivoting does force the entire current
tail to vanish. The lemma requires neither a zero-target extension
nor equality of the two full-matrix constructions at a zero-product
junction.

Actual sign changes on the final rows and columns reduce every
\(X\)-tail to
\begin{equation}\label{eq:Xtail}
H=\begin{pmatrix}r&r\\r&w\end{pmatrix},
\qquad r>0,\quad |w|\le r,\quad F=r-w.
\end{equation}
Here \(H\) is the \emph{actual} fourth-stage tail, whose entries are
\(k\) times the normalized core-tail entries.
Indeed, for
\((\begin{smallmatrix}\epsilon_1r&\epsilon_2r\\
\epsilon_3r&W\end{smallmatrix})\),
multiply rows by \((\epsilon_1,\epsilon_3)\) and columns by
\((1,\epsilon_1\epsilon_2)\). This gives \eqref{eq:Xtail}, with
\(w=\epsilon_1\epsilon_2\epsilon_3W\), and preserves all absolute
matrix bounds.
Similarly every \(D\)-tail has a form
\begin{equation}\label{eq:Dtail}
H=\begin{pmatrix}r&s\\t&\chi r\end{pmatrix},
\quad r>0,\quad0\le s,t\le r,\quad\chi\in\{-1,1\}.
\end{equation}
If \(\chi=1\), its height is \(r-st/r\le r\le4\); a zero arm
also gives height \(r\le4\).
If \(s=r\) or \(t=r\), exchanging the corresponding tied final
columns or rows and normalizing signs gives \eqref{eq:Xtail} at the
same height. The sole additional domain is therefore
\begin{equation}\label{eq:properB}
H=\begin{pmatrix}r&s\\t&-r\end{pmatrix},
\qquad0<s,t<r,\qquad F=r+\frac{st}{r}.
\end{equation}
We call this the proper negative diagonal domain.

\section{The complete cross-saturated bound}\label{sec:X}

For \eqref{eq:Xtail}, the actual core matrix is
\begin{equation}\label{eq:XD}
D=\begin{pmatrix}
k&A&B\\
ck&r+cA&r+cB\\
dk&r+dA&w+dB
\end{pmatrix}.
\end{equation}
All of \eqref{eq:bands1}--\eqref{eq:bands4} and \(\mx D\le k\) are
retained. No rank, cofactor-sign, or auxiliary entry condition is assumed.

\begin{proposition}[Certified complete \(X\) bound]\label{prop:X}
Every real complete lift with \eqref{eq:XD}, \(p,k,r>0\), and
\(|w|\le r\), satisfies \(F=r-w\le\alpha\).
\end{proposition}

The proposition is assembled from four groups of certified source domains.
Appendix~\ref{app:artifacts} and Supplementary Material~S4
identify the component certificates and their source domains.
\begin{center}
\begin{tabular}{@{}p{.53\linewidth}p{.33\linewidth}@{}}
\toprule
Complete source domain&Certified bound\\
\midrule
\(0<k\le2\)&\(F<1653/400\)\\
\(k\ge21/10\)&\(F<\gam\)\\
\(2<k\le21/10,\ r>k\)&\(F<\gam\)\\
\(\gam/2\le k\le21/10,\ r\le k,\ F\ge\gam\)&\(F\le\alpha\)\\
\bottomrule
\end{tabular}
\end{center}

\begin{proof}[Assembly]
The first three rows leave \(2<k<21/10,\ r\le k\).
If \(F<\gam\), \eqref{eq:thresholds} suffices. Otherwise
\(\gam\le F\le2r\le2k\), hence \(k\ge\gam/2\), and the last row
applies. The boundary \(r=k\) lies in the last row; \(k=2,21/10\)
have closed coverage. Internal sign and transpose normalizations are
actual matrix operations imposing no extra rank or cofactor hypothesis.
\end{proof}

The small-\(k\) certificate has \(268817\) nodes, \(134409\)
contradiction leaves, and no open leaf. The accepted two low-\(r\)
roots have the following counts:
\begin{center}
\begin{tabular}{@{}lrrrr@{}}
\toprule
Root&Nodes&Contradiction leaves&\(\alpha\)-safe leaves&Open\\
\midrule
I&3450095&1725047&1&0\\
II&5920895&2873490&86958&0\\
\bottomrule
\end{tabular}
\end{center}
The counts identify accepted objects, rather than replacing their
acceptance rules and domain maps. An \(\alpha\)-safe leaf can contain
feasible points; it is not an infeasibility certificate at \(\gam\).

\paragraph{Order of the analytic inputs.}
The low-\(r\) certificate uses a wall flow in the complete
22-coordinate \(X\) model. This flow decreases \(r+w\) to zero;
only at that wall is a signed tail permutation applied. Its endpoint
satisfies all sixteen strict positive-room conditions, and its height
is bounded by the independent strict-room theorem. That theorem
uses the two-face normalization, the safe-tail and middle-face
reductions, and the canonical-height theorem explained below.
The separate weak-room exit uses the full six-cofactor closure
theorem, including its zero-coordinate cases.
In contrast, the two-gap transport of Lemma~\ref{transport:local-lemma}
uses 24 source coordinates, decreases \(r-s\) and \(r-t\),
and preserves \(r+w\). Its endpoint is bounded by
Proposition~\ref{prop:X}. Neither that transport nor the later
guard-preserving flow is an input to the low-\(r\) certificate or to
Proposition~\ref{prop:X}. The distinct original statements and
their entry conditions are identified in Supplement S4.

% Included section: s3_main
\subsection{The exact height in the canonical near-peak domain}
\label{s3:main}

The near-peak part of the argument requires an algebraic height bound.
A terminal estimate of the form \(F\le\amin<\alpha\) cannot cover a
matrix whose height is exactly \(\alpha\). We therefore explain the
separate argument which supplies the exact bound, before using it in
local transports and in the complete \(X\)-domain certificate.

Let \(\mathcal K\) be the positive canonical chamber in
Definition~\ref{s3:canonical-domain}. Its matrices have the complete
physical representation of Section~\ref{sec:lift}, with
\(e=1,\ u=(1,V,T)^T\), \(P_1=-1\), \(0<V<1\), the stated positive receiver
coordinates and six strictly positive cofactors. All initial,
intermediate, and core entry bounds remain part of the definition.
No tail contact, equality of row scales, or equality-family hypothesis
is added to this chamber.

\begin{theorem}[Canonical exact height]\label{s3:canonical-height-main}
Every complete physical matrix in \(\mathcal K\) satisfies
\(F\le\alpha\). More precisely, every admissible target representative
of Definition~\ref{s3:target-representative} has height at most
\(\alpha\). The representative and reconstruction results in
Propositions~\ref{s3:target-reduction} and~\ref{s3:scale-equivalence}
then give the assertion for every matrix in \(\mathcal K\).
\end{theorem}

\begin{proof}
We give the argument here; Appendix~\ref{s3:details} specifies the
complete source predicate, reconstruction formulas, eight constraints,
and finite verification statements used in the proof.
Set
\[
 f_0=\frac{4132517}{10^6},\qquad
 f_1=\frac{103313}{25000},\qquad f_0<\alpha<f_1.
\]
Suppose a source has height greater than \(\alpha\). The target
reduction provides, at the same seven-coordinate head, a complete
representative of any chosen height \(f\) between \(\alpha\) and the
source height, with \(f<5\). It does not assert that the source itself
already has the representative's contacts.

In the row-scale coordinates write
\[
 p=1+s,\qquad r=1+R,\qquad F=2(1+R),
 \qquad z=(\beta,V,s,Y,Z,v,h,b,R).
\]
The variables \(t,E,X,k,g,j,a\) and the complete matrix are recovered
from this one tuple by \eqref{s3:scale-dictionary}.
The equivalence in Proposition~\ref{s3:scale-equivalence} retains
both column capacities and the strict scale qualifications; it is
not a relaxation obtained by completing different columns separately.

First, a finite rational certificate puts every height-\(f_0\)
representative into the six-dimensional box \(Q\) in
\eqref{s3:Q}. Comparing target representatives at the same old head
is legitimate because their tail receiver satisfies the identity
\[
 v(R)=\frac{R+Y-t(1+X)}{t(1-\beta X)-\Gamma},
 \qquad \Gamma=V-\beta Y.
\]
On the height-\(f_0\) representatives the denominator is greater
than \(1/5\). Consequently, increasing the target from \(f_0\) to
\(f_1\) changes \(v\) by less than \(10^{-5}\), while the other
five localization coordinates stay fixed. A separate finite certificate
excludes height \(f_1\) in this enlarged box. Every potentially
supercritical representative therefore has \(F<f_1\) and lies in the
nine-dimensional box \(\mathcal B\) of \eqref{s3:nine-box}.
This step is what connects a fixed-target localization certificate to
an upper bound for all possible source heights.

Second, let \(G_1,\ldots,G_8\) be the physical slacks in
\eqref{s3:eight-slacks}. On \(\mathcal B\), all the other strict
qualifications have uniform positive margins. The exact
reconstruction theorem thus identifies the high-value completion set
with
\[
 \{z\in\mathcal B:G_1(z),\ldots,G_8(z)\ge0\}.
\]
It is compact. Under the assumed counterexample, its maximum height
is greater than \(\alpha\); no separate numerical witness is needed
for this existence argument.
For fixed \(s\), the Jacobian
\[
 J=\frac{\partial(G_1,\ldots,G_8)}
 {\partial(\beta,V,Y,Z,v,h,b,R)}
\]
is invertible throughout \(\mathcal B\). The uniform rational
certificate in Proposition~\ref{s3:jacobian} gives
\[
 \frac{\partial R}{\partial G_i}<-\frac1{250}
 \quad(1\le i\le8).
\]
If a maximizing point had \(G_i>0\), the local inverse-function
coordinates would allow that slack to decrease, with \(s\) and all
other slacks fixed. This would preserve the physical inequalities
and strictly increase the height. A short such arc is an arc of the
original completion model; the proved global localization prevents
an artificial box boundary from invalidating the argument.
Hence a maximizing point satisfies all eight equalities.

Third, these equalities force
\[
\begin{gathered}
 t=X=Z=1,\quad h=T=s,\quad b=k=r+Ea,\\
 O_{01}=-1,\quad S_{01}=-p,\quad O_{21}=1.
\end{gathered}
\]
Only at this stage is the contact family used.
With \(\epsilon=r-2\), its remaining equation is the explicitly
defined polynomial \(G(s,\epsilon)=0\) in \eqref{s3:carrier-poly}.
The same invertible eight-constraint Jacobian makes the contact
family a local smooth curve parametrized by \(s\).
At its maximum, \(\epsilon'(s)=0\), so differentiation gives
\(G_s(s,\epsilon)=0\). This conclusion does not require a separate
assumption that \(G_\epsilon\ne0\).

Finally, exact elimination gives
\[
 \operatorname{Res}_s(G,G_s)=
 \frac{\epsilon^{34}(\epsilon-2)(\epsilon-1)^2(\epsilon+1)^7
       (\epsilon^2-2\epsilon+2)^2}{131072}
 P_{61}(4+2\epsilon).
\]
Every factor outside \(P_{61}\) is nonzero in the proved height
interval. By Section~\ref{sec:alpha}, the polynomial \(P_{61}\) has
exactly one root in \((4,5)\), namely \(\alpha\), and therefore
exactly one in \([f_0,f_1]\). Thus the maximum must equal \(\alpha\), contrary to the
assumed larger height.
\end{proof}

The proof uses a constrained maximum argument before it invokes
stationarity and root identification. It therefore does not infer a
neighborhood bound from an isolated critical point, nor does it use
the restricted-family optimum in
Chen--Edelman--Urschel~\cite{ChenEdelmanUrschel2026}, Theorem~2.2,
as a substitute for reaching that family. The source reductions and
the finite certificates have distinct roles, all specified in
Appendix~\ref{s3:details}. The local \(X\)-transports first reach
an actual positive-room matrix, to which this canonical result and
the explicit room-reduction results apply. The complete \(X\)
bound subsequently serves the negative-diagonal proof; it is not an
input to the canonical-height argument.


\section{Exact reduction of the negative diagonal domain}\label{sec:fibre}

Replace \eqref{eq:XD} by
\begin{equation}\label{eq:BD}
D=
\begin{pmatrix}
k&A&B\\
ck&r+cA&s+cB\\
dk&t+dA&-r+dB
\end{pmatrix}.
\end{equation}
We initially allow \(0\le s,t\le r\), so the safe boundary is included.
A \emph{complete source} means \eqref{eq:lift},
\eqref{eq:bands1}--\eqref{eq:bands4}, \(\mx D\le k\), and these tail
conditions.
The normalized first Schur tail is a real \(4\times4\) completely
pivoted matrix. Corollary~\ref{cor:firstfour} gives
\begin{equation}\label{eq:F4p}
F\le4p.
\end{equation}
Hence \(F>4\) requires \(p>1\). The head bound implies
\(e\beta\ge p-1>0\). A simultaneous sign flip of the second row and
column gives \(e,\beta>0\). Simultaneously reversing all three remaining
rows and columns leaves \(D\) fixed and reverses \(x\), allowing
\(x_0\ge0\). Simultaneously reversing just the last two rows and
columns reverses \(A,B,c,d\), leaves \(x_0\) fixed, and leaves the
last tail fixed, allowing \(A\le0\).
No sign is imposed on the other
receivers or on \(c,d\).

Fix the seventeen-coordinate common framework
\begin{equation}\label{eq:frame}
\vartheta=(k,A,B,c,d,u_0,u_1,u_2,x_0,x_1,x_2,
v_0,v_1,v_2,q_0,q_1,q_2).
\end{equation}
The purpose of keeping this common framework fixed is to separate two
questions: which tail values are attainable by one complete matrix, and
how large the best such tail can be. The following three steps answer
both questions exactly. In contrast to independent entry bounds, the
resulting capacities are realized simultaneously.

We maximize over the six remaining actual variables
\((\beta,p,e,r,s,t)\), in the normalized range
\(p\ge1\), \(0\le e,\beta\le1\). Every high source has such a
representative by the preceding argument.
The constructions give feasible paths, not
independent necessary relaxations.

\subsection{The right-prefix interval}
Define
\[
I_\beta=[0,1]\cap\bigcap_{j=0}^2\{b:|q_j+bv_j|\le1\}.
\]
If \(v_j=0\), this condition is \(|q_j|\le1\); otherwise its endpoints
are \((-1-q_j)/v_j\) and \((1-q_j)/v_j\), sorted in increasing order.
An empty interval rules out the fibre. When it is nonempty, set
\(\bar\beta=\max I_\beta\).

\begin{lemma}\label{lem:beta}
For a complete source with \(p\ge1\), \(0\le e,\beta\le1\),
increasing \(\beta\) linearly to \(\bar\beta\), while fixing the
other variables, preserves complete feasibility and \(F\).
\end{lemma}
\begin{proof}
The right-prefix bands define the convex interval \(I_\beta\).
The head upper bound \(p-e\beta\le1\) improves, and its lower bound is
automatic because \(p\ge1\) and \(e\beta\le1\).
No other condition or height depends on \(\beta\).
\end{proof}

\subsection{The attained \texorpdfstring{\(p,e\)}{p,e} capacity}
If \(\bar\beta=0\), the head forces \(p=1\), safe by \eqref{eq:F4p}.
Suppose \(\bar\beta>0\). Form the head row
\((a_0,b_0)=(1,\bar\beta)\). For every \(x_i\ne0\), add the row
\[
a_i=|x_i|>0,\qquad b_i=\sgn(x_i)u_i.
\]
When \(x_i=0\), its prefix bound is automatic.
There are at most four rows, with \(0<a_i\le1\), \(|b_i|\le1\).
The complete remaining prefix system is
\begin{equation}\label{eq:prefixpolygon}
a_ip-b_ie\le1,\qquad0\le e\le1,\qquad p\ge1.
\end{equation}
The reverse bound is automatic from \(a_ip\ge0\) and \(b_ie\le1\).
Define capacities for \(b_i\ge0\) and for \(b_i>0>b_j\), respectively:
\[
C_i=\frac{1+b_i}{a_i},\qquad
C_{ij}=\frac{b_i-b_j}{a_jb_i-a_ib_j}.
\]
All denominators are positive. A single row bounds \(p\) by using
\(e\le1\); a positive and a negative row bound it by requiring their
lower and upper limits for \(e\) to intersect. There are no other
obstructions in this two-dimensional polygon. Put
\begin{equation}\label{eq:Pestar}
P=\min(C_i,C_{ij}),\qquad
e_*=\max\left(0,\max_{b_i>0}\frac{a_iP-1}{b_i}\right).
\end{equation}

\begin{lemma}\label{lem:prefixcap}
The maximum \(p\) in \eqref{eq:prefixpolygon} is \(P\);
\((P,e_*)\) is feasible. Furthermore \(1\le P\le2\), and at most six
capacities occur.
\end{lemma}
\begin{proof}
Writing \(\ell_i=1-a_ip+b_ie\ge0\), feasibility implies
\[
1+b_i-a_ip=\ell_i+b_i(1-e)\ge0\quad(b_i\ge0),
\]
and
\[
b_i-b_j-(a_jb_i-a_ib_j)p=b_i\ell_j-b_j\ell_i\ge0
\quad(b_i>0>b_j).
\]
Thus every capacity is necessary.
Conversely positive \(b_i\) give lower bounds for \(e\), negative
\(b_j\) give upper bounds, and zero \(b_i\) constrain only \(p\).
The \(C_i\) make all lower bounds at most one; the \(C_{ij}\) make
each lower bound at most each upper bound.
A negative row additionally requires a nonnegative upper bound,
namely \(p\le1/a_j\). Its pairing with the head implies this:
\[
\frac{\bar\beta-b_j}{a_j\bar\beta-b_j}\le\frac1{a_j},
\]
since \(a_j\le1\), \(b_j<0\).
The head lower bound \((P-1)/\bar\beta\) is nonnegative.
Therefore \eqref{eq:Pestar} satisfies every row.
Each capacity is at least one, and the head capacity is
\(1+\bar\beta\le2\). If the numbers of positive, zero, and negative
rows are \(m,z,n\), then their number of capacities is \(m+z+mn\),
where \(m\ge1\), \(m+z+n\le4\), and hence at most six.
\end{proof}

After Lemma~\ref{lem:beta}, move \((p,e)\) along the segment to
\((P,e_*)\). The prefix polygon is convex and \(p\) increases.
The actual matrices
\[
D,\qquad S=D+xq^T,\qquad O=D+xq^T+uv^T
\]
stay fixed. Thus \(|q_j|\le p\), \(|S_{ij}|\le p\) improve, and
all other bands remain valid. This second path preserves \(F\);
it does not require \(e\) itself to vary monotonically.

\subsection{The attained tail maximum}
At \(p=P\), form joint intervals for the entries of \(D\):
\begin{equation}\label{eq:entrycaps}
\begin{split}
\ell_{ij}&=\max(-k,-P-x_iq_j,-1-x_iq_j-u_iv_j),\\
U_{ij}&=\min(k,P-x_iq_j,1-x_iq_j-u_iv_j).
\end{split}
\end{equation}
The five entries in the zeroth row or column of \(D\) are fixed.
Check their intervals and \(|q_j|\le P\). Failure makes the fibre
empty: decreasing \(p\) cannot repair a failed stage bound, while
original-matrix and core bounds do not depend on \(p\).
For the other entries put
\begin{align*}
r_-&=\max(0,\ell_{11}-cA,dB-U_{22}),
&R&=\min(U_{11}-cA,dB-\ell_{22}),\\
s_-&=\max(0,\ell_{12}-cB),
&s_+&=U_{12}-cB,\\
t_-&=\max(0,\ell_{21}-dA),
&t_+&=U_{21}-dA,\\
S_*&=\min(R,s_+),
&T_*&=\min(R,t_+).
\end{align*}

\begin{theorem}[Complete six-variable fibre maximum]\label{thm:fibre}
For a nonempty high-value fibre at fixed \(\vartheta\), the
construction yields a complete source with
\[
(\beta,p,e,r,s,t)=(\bar\beta,P,e_*,R,S_*,T_*),
\]
and
\begin{equation}\label{eq:Bcap}
\cB(\vartheta)=R+\frac{S_*T_*}{R}
=\max_{\text{complete six-variable fibre}} F.
\end{equation}
The tail fibre is nonempty exactly when the fixed checks pass and
\[
R>0,\qquad R\ge r_-,\qquad S_*\ge s_-,\qquad T_*\ge t_-.
\]
Every original source reaches the maximum by three complete feasible
paths with nondecreasing height.
\end{theorem}

\begin{proof}
The first two paths have been established. At their endpoint the
tail feasible set is exactly
\[
r_-\le r\le R,\quad s_-\le s\le s_+,\quad
t_-\le t\le t_+,\quad r>0,\quad0\le s,t\le r.
\]
The intervals are computed simultaneously from the same \(u,x,v,q\).
Nonemptiness implies \(R\ge r>0\), \(S_*\ge s\), \(T_*\ge t\);
then \((R,S_*,T_*)\) satisfies every interval. The stated conditions
are therefore necessary and sufficient. The endpoint dominates every
feasible tail coordinatewise. Coordinatewise domination would not suffice
for an arbitrary nonlinear objective: here the complete-pivoting condition
\(st\le r^2\) supplies the needed sign. The following identity makes
that dependence explicit. For any feasible tail,
\begin{align}
Rr(\cB-F)
={}&(R-r)(Rr-st)+rt(S_*-s)+rS_*(T_*-t)\ge0.\label{eq:taildifference}
\end{align}
The factors have the required signs because \(R\ge r\) and
\(st\le r^2\). The segment to this maximum stays in the convex
feasible set, with all three coordinates nondecreasing. Along it,
\[
\partial_rF=1-st/r^2\ge0,\qquad
\partial_sF=t/r\ge0,\qquad \partial_tF=s/r\ge0.
\]
This proves the last path and the attained maximum.
\end{proof}

Every hypothetical \(F>\alpha\) source has its framework in the closed
box
\begin{equation}\label{eq:Broot}
\begin{gathered}
\frac{4\gam}{9}\le k\le\frac94,\quad
-\frac94\le A\le0,\quad-\frac94\le B\le\frac94,\quad
|c|,|d|\le1,\\
|u_i|,|v_i|\le1,\quad0\le x_0\le1,\quad
|x_1|,|x_2|\le1,\quad |q_i|\le2.
\end{gathered}
\end{equation}
The lower bound on \(k\) follows from \(F\le(9/4)k\), and the other
bounds from the complete bands and actual sign normalizations.
The root retains the complete fibre conditions, not merely the box.
In particular it does not discard the proper negative domain with
\(k\le2\). An additional shared head budget \(pk\le4\), when used,
follows from the leading \(3\times3\) determinant: its magnitude is
\(pk\), and the maximum determinant on the real unit-entry cube is
\(4\). For completeness, multilinearity reduces that maximum to
sign matrices; their \(3\times3\) determinants are multiples of
four, while the Hadamard bound is \(3\sqrt3<8\).

All intersections, minima, and maxima include ties.
Zero coefficients are handled before division. The point \(R=0\)
cannot support a high source and is rejected as an empty
positive-tail fibre. No early generic-position restriction is needed.

% Included section: root_semantics
\subsection{The complete source carried by a certificate}\label{sec:source-interface}
It is useful to state explicitly what the root box represents. A point
of the box alone need not admit a complete matrix, and an auxiliary
product variable alone need not equal a product of physical coordinates.
The following definition fixes the mathematical source throughout the
computational argument.

\begin{definition}[Complete fibre and high framework]\label{def:complete-fibre}
For \(\vartheta\) in \eqref{eq:frame}, let
\(\mathcal F(\vartheta)\) be the set of tuples
\(\zeta=(\beta,p,e,r,s,t)\) such that
\[
p\ge1,\quad k,r>0,\quad 0\le e,\beta\le1,
\quad 0\le s,t\le r,
\]
and the matrix defined by \eqref{eq:lift} and \eqref{eq:BD} satisfies
all of \eqref{eq:bands1}--\eqref{eq:bands4} and \(\mx D\le k\).
Let \(Q_0\) denote the closed box \eqref{eq:Broot}. Define
\[
\mathcal H_\gamma=
\left\{\vartheta\in Q_0:
\exists\zeta\in\mathcal F(\vartheta),\quad
r+\frac{st}{r}\ge\gamma\right\}.
\]
All inequalities in this definition refer to the same reconstructed
matrix. Zero tails, whose height is zero, do not belong to this high domain.
\end{definition}

\begin{proposition}[Root semantics]\label{prop:root-semantics}
Every normalized proper negative-diagonal source with \(F>\alpha\)
has, after the actual sign changes above, a framework in
\(\mathcal H_\gamma\). For every \(\vartheta\in\mathcal H_\gamma\),
Theorem~\ref{thm:fibre} constructs a member
\(\zeta_*(\vartheta)\in\mathcal F(\vartheta)\) with
\[
F(\vartheta,\zeta_*(\vartheta))=\cB(\vartheta)
=\max_{\zeta\in\mathcal F(\vartheta)}F(\vartheta,\zeta).
\]
Thus the negative-diagonal computation must establish
\(\cB(\vartheta)\le\alpha\) for all
\(\vartheta\in\mathcal H_\gamma\).
\end{proposition}
\begin{proof}
The sign normalizations preserve height and feasibility. The bound
\(F\le4p\) and \(\alpha>4\) give \(p>1\), and
\eqref{eq:Broot} contains the normalized framework. Since
\(F>\alpha>\gamma\), this framework belongs to
\(\mathcal H_\gamma\). The attained endpoint and equality with the
fibre maximum follow from Theorem~\ref{thm:fibre}. All divisions and
selector choices in that theorem are qualified by its nonempty
high-fibre hypothesis.
\end{proof}

The explicit interval construction is also a decision procedure for
these qualifications: it checks the right-prefix interval, the prefix
capacities, the five fixed core entries, and the simultaneous tail
intervals. Empty intersections are rejected; equal endpoints and tied
minima are included. A zero coefficient is handled by its original
inequality before any division. If a formula requires a sign that is
not established over a box, its alternatives must remain covered.
These cases are part of the source semantics, rather than discretionary
choices of the search algorithm.

For clarity, a necessary polynomial model may describe a larger set
than \(\mathcal H_\gamma\). This is legitimate only after proving that
every relevant complete source supplies all its variables
simultaneously and satisfies every model row. Excluding the larger set
then excludes that source. A feasible-transport terminal requires more:
its image must be an actual complete source, because a bound for a
relaxed auxiliary point does not by itself bound a physical matrix.

\begin{center}
\small
\begin{tabular}{@{}>{\raggedright\arraybackslash}p{.24\linewidth}>{\raggedright\arraybackslash}p{.67\linewidth}@{}}
\toprule
Operation&Required mathematical relation\\
\midrule
Introduce products&One complete point supplies each shared product with its exact value.\\
Form necessary rows&Every high source in the current region satisfies those same rows.\\
Intersect a box&Every source lost by the intersection is already impossible or safe.\\
Select a capacity formula&Its signs and selector conditions hold, or all alternatives are retained.\\
Change coordinates&The same physical source is reconstructed, or a complete image of no smaller height is supplied.\\
Attach a subproof&Its full input source, including the inherited transformation, agrees with the residual obligation.\\
\bottomrule
\end{tabular}
\end{center}


\section{Rational certificates and mixed-height coverage}
\label{sec:certificates}

The remaining finite proof obligation on \eqref{eq:Broot} is that every
complete framework has \(\cB(\vartheta)\le\alpha\).
Discovery proposes subdivisions and terminal certificates.
Acceptance checks them by exact rational arithmetic on their
complete declared sources. Floating-point optimization may guide
discovery but is not an acceptance criterion.

\subsection{Necessary outer systems}
Each actual source gives all coordinates and products simultaneously.
Necessary polynomial inequalities may be relaxed over a rational box,
provided the relaxation contains this same source. For example, for
\(L_x\le x\le U_x\), \(L_y\le y\le U_y\), the shared product \(z=xy\)
satisfies
\begin{align*}
z&\ge L_xy+L_yx-L_xL_y,&
z&\ge U_xy+U_yx-U_xU_y,\\
z&\le L_xy+U_yx-L_xU_y,&
z&\le U_xy+L_yx-U_xL_y.
\end{align*}
A product appearing in several rows is not assigned independent
physical values. Enclosures are outward. Sign-sensitive division or
use of a selected min/max formula requires qualification on the entire
box, or subdivision that retains every alternative.

If the necessary rows are \(a_j\cdot z\le b_j\), exact nonnegative
weights give
\[
a=\sum_j\lambda_ja_j,\quad
B=\sum_j\lambda_jb_j,\quad
a\cdot z\le B,\quad\lambda_j\ge0.
\]
A box lower bound for \(a\cdot z\) greater than \(B\) rules out the
source. The acceptor checks all weights, row identities, endpoints,
and strict comparisons as rationals.
An interval contradiction is the one-expression case.

These are inequalities for arbitrary real points in the box. The
rational endpoints and weights make the finite arithmetic exact; they
do not restrict the matrix entries to rational numbers. Thus a verified
strict separation excludes a continuum of real sources at once. The finite
computation is a proof of a universal implication on a region, rather
than a sampling of matrices from that region.
The actual height can be introduced as a shared coordinate through
\[
rF-r^2-st=0,
\]
using both inequality directions. It must not be replaced by the
\(X\)-specific objective \(r-w\).

% Included section: certificate_example
% Insert within Section 8, after the necessary-row acceptance rule.
\subsection{A real source-bound contradiction certificate}
\label{certificate:worked}

We give an actual terminal from the final negative-diagonal computation.
It lies at the local binary path
\(\mathtt{1001110110011011111}\) in parent 435003, below the anchor
\(\mathtt{10011101}\). The source being excluded consists of the
canonical high sources at this target satisfying \(F\ge\gam\), with
the complete parent and ancestor restrictions retained. It is not the
entire geometric box. The frozen construction supplies 24 rationally
bounded source coordinates and 43 shared products, hence a lifted vector
\(z\in\R^{67}\), together with 692 necessary linear inequalities.
Every physical source under consideration has a lift satisfying these
inequalities and all 67 coordinate bounds.

The certificate selects 49 rows with positive integer weights; all
unselected weights are zero. Two of its actual rows are
\begin{center}
\begin{tabular}{@{}cl r@{}}
\toprule
Row index & Necessary inequality & Weight \\
\midrule
2 & \(-z_0-z_4\le0\) & \(31\,317\,836\) \\
9 & \(-z_0+z_2+z_{36}\le0\) & \(229\,060\,761\) \\
\bottomrule
\end{tabular}
\end{center}
Indices here are zero-based. The first 24 lifted coordinates are the
source coordinates; in this dictionary \(z_{36}=z_5z_6\).
The two displayed rows are an extract, not a two-row certificate:
the following calculation uses all 49 selected rows.

Write their exact weighted sum as \(a\cdot z\le B\). If
\([L_i,U_i]\) is the bound for coordinate \(i\), then
\begin{equation}\label{certificate:boxminimum}
m=\sum_{i=0}^{66}
\begin{cases}
a_iL_i,&a_i\ge0,\\
a_iU_i,&a_i<0
\end{cases}
\quad\text{satisfies}\quad m\le a\cdot z.
\end{equation}
For example, the actual first coordinate has
\[
\begin{gathered}
[L_0,U_0]=
\left[\frac{142554217571}{68719476736},
      \frac{142566885919}{68719476736}\right],\\
a_0=-\frac{237458354877801951091349245551}
                {1267650600228229401496703205376}.
\end{gathered}
\]
Thus its contribution to \(m\) uses \(U_0\). Summing the 67
contributions and subtracting the exact weighted right-hand side gives
\begin{equation}\label{certificate:actualmargin}
m-B=
\frac{\scriptstyle
1005495384602754603811064507351680469029685784938386479}
{\scriptstyle
544451787073501541541399371890829138329600000000000}
>0.
\end{equation}
Consequently \(m\le a\cdot z\le B<m\) is impossible. This is an
exact rational contradiction for the specified high source. Its margin
depends on the normalization of the weights; it is not a bound on
\(\alpha-F\).

The complete row list, weights, coordinate bounds and weighted sums
are supplied in the real example identified in
Appendix~\ref{certificate:details}. Its source record also identifies
the frozen parent, projection, ancestor trace and target image. These
bindings are needed to establish that the rows and bounds apply to
this source; the arithmetic in
\eqref{certificate:actualmargin} alone does not establish that fact.

This terminal pays one position in the anchor's complete 54-terminal
cover. The other positions are retained, including those proved safe
by a height bound or feasible transport instead of a contradiction.
By Lemma~\ref{lem:mixed}, their composition proves that complete anchor
safe at \(\alpha\). The anchor then pays the last of 240 constituent
sources of parent 435003; the other 239 sources retain their accepted
coverage. Finally, that parent is bound to one of the four open
positions of the preserved original root. Thus the example illustrates
both parts of certificate acceptance: an exact local inequality and a
complete, source-preserving route back to the original domain.


\subsection{Height certificates and feasible transports}
Suppose \(F=h\cdot z\), after introducing a height coordinate if
necessary. For \(t>0\), let \(c=a-th\) and
\[
m=\sum_i\min(c_iL_i,c_iU_i).
\]
Then
\begin{equation}\label{eq:heightcert}
tF=a\cdot z-c\cdot z\le B-m.
\end{equation}
Thus \(B-m\le t\amin\) proves \(F\le\alpha\).
It does not establish infeasibility.

Other safe terminals use feasible transformations to a domain already
bounded by \(\alpha\). If every source \(z\) in a box has a complete
image \(T(z)\) with \(F(T(z))\ge F(z)\), and
\(F(T(z))\le\alpha\), that box is safe.
The transformation must preserve the actual coupled matrix bands,
rather than changing unrelated abstract capacities.

% Included section: transport_main
The following local transport gives a quantitative safe terminal near
cross-saturated tails.  Its parameter bounds hold on full-dimensional
boxes; the construction is not restricted to an equality family.
Write the complete source coordinates as
\begin{equation}\label{transport:coordinates-main}
\begin{split}
\xi={}&(k,r,w,A,B,c,d,p,e,\beta,u_0,u_1,u_2,x_0,x_1,x_2,\\
&\hspace{28mm}v_0,v_1,v_2,q_0,q_1,q_2,\sigma,\tau)\in\R^{24},
\end{split}
\end{equation}
where \(\sigma=r-s\) and \(\tau=r-t\).  Let \(y\in\R^{23}\)
be the ordered tuple obtained by deleting \(p\) from \(\xi\).
Thus \(\xi\) is specified by \((p,y)\), with the ordering in
\eqref{transport:coordinates-main} understood.  The actual tail and
its height expression are
\begin{equation}\label{transport:tail-height}
H(\xi)=\begin{pmatrix}r&r-\sigma\\r-\tau&w\end{pmatrix},
\qquad F(\xi)=\frac{(r-\sigma)(r-\tau)}r-w.
\end{equation}
A complete physical source satisfies the original entry and successive
Schur-complement bounds, as well as
\(\sigma,\tau,r-\sigma,r-\tau\ge0\).  Appendix~\ref{transport:details}
specifies these bounds as 104 polynomial slacks and gives their actual
matrix reconstruction.  In particular, the proper negative-diagonal
domain \(w=-r\), \(\sigma,\tau>0\) is a subdomain of the sources
considered below; it is not an additional assumption of the transport.

For each \(j\in\{0,1,2,3\}\), select the twenty physical slacks
\(\ell_{j,0},\ldots,\ell_{j,19}\) listed in
Table~\ref{transport:slack-table}, and set
\begin{equation}\label{transport:coordinate-map}
\Phi_{j,p}(y)=
 (\ell_{j,0}(\xi),\ldots,\ell_{j,19}(\xi),r+w,\sigma,\tau),
\qquad J_j(\xi)=D_y\Phi_{j,p}(y).
\end{equation}
Holding \(p\) fixed in this derivative does not restrict its allowed
value to a centre coordinate.  All four rational centres \(c_j\) have
the same \(p\)-coordinate
\[
p_* = \frac{726612331}{500000000}.
\]
The estimates below hold on each entire closed box
\begin{equation}\label{transport:outer-box}
Q_j=\{\xi:\iin{\xi-c_j}\le R\},\qquad R=\frac1{1250},
\end{equation}
and hence uniformly for the full parameter interval
\begin{equation}\label{transport:p-interval}
p\in\left[\frac{726212331}{500000000},
            \frac{727012331}{500000000}\right].
\end{equation}
The rational centres, the selected slacks, and the finite verification
identities are specified in Appendix~\ref{transport:details}.

\begin{lemma}[Uniform local transport to a complete \(X\) source]
\label{transport:local-lemma}
Let \(\xi_0\) be a complete physical source, with gaps
\(\sigma_0,\tau_0\ge0\).  Suppose that, for one of the four centres,
\begin{equation}\label{transport:strict-start}
d_0+3(\sigma_0+\tau_0)<R,
\qquad d_0=\iin{\xi_0-c_j}.
\end{equation}
There is a piecewise continuously differentiable path of complete
physical sources from \(\xi_0\) to a complete \(X\) source \(\xi_*\)
inside \(Q_j\).  Along this path, \(p\), \(r+w\), and the twenty
selected slacks retain their values at \(\xi_0\).  Moreover,
\[
F(\xi_*)-F(\xi_0)\ge\frac{\sigma_0+\tau_0}{10}.
\]
Consequently, Proposition~\ref{prop:X} gives
\begin{equation}\label{eq:gapsafe}
F(\xi_0)\le\alpha-\frac{\sigma_0+\tau_0}{10}.
\end{equation}
The selected slacks need not vanish at the source.
\end{lemma}

\begin{proof}
We give the uniform estimates and the continuation argument explicitly.
For each \(j\), the rational preconditioner \(C_j\) satisfies, on all
of \(Q_j\),
\begin{equation}\label{transport:contraction-main}
\iin{I-C_jJ_j(\xi)}\le
\eta:=\frac{396725479558446119}{10^{19}}<\frac1{25}.
\end{equation}
Here the matrix norm is the induced infinity norm.  Thus \(J_j(\xi)\)
is invertible.  If \(e_\sigma,e_\tau\) denote the last two coordinate
vectors in the range of \(\Phi_{j,p}\), put
\[
v_\sigma(\xi)=J_j(\xi)^{-1}e_\sigma,
\qquad v_\tau(\xi)=J_j(\xi)^{-1}e_\tau.
\]
The same whole-box rational estimates give
\begin{equation}\label{transport:direction-main}
\begin{gathered}
\iin{v_\sigma},\ \iin{v_\tau}\le
\frac{24566153053000000000}{9603274520441553881}<3,\\
D_yF\,v_\sigma\le-\frac1{10},\qquad
D_yF\,v_\tau\le-\frac1{10}.
\end{gathered}
\end{equation}
Every one of the 81 slacks not used as a component of
\(\Phi_{j,p}\) is bounded below throughout \(Q_j\) by
\begin{equation}\label{transport:margin-main}
\mu:=\frac{913752034886059}{10^{18}}>0.
\end{equation}
In particular the first four pivots \(1,p,k,r\) stay nonzero.
Also, throughout these boxes,
\[
r\ge r_{\min}:=\frac{2065458539}{10^9}>0,
\qquad w\le-r_{\min}<0.
\]
Together with the physical inequalities \(r-\sigma,r-\tau\ge0\),
this makes \(F\) in \eqref{transport:tail-height} positive, so that
it is the magnitude of the final pivot of the reconstructed matrix.
Appendix~\ref{transport:details} derives
\eqref{transport:contraction-main}--\eqref{transport:margin-main}
from explicit polynomial coefficients and rational matrices, without
using a numerical trajectory.

Fix \(p=p(\xi_0)\).  First solve the autonomous equation
\begin{equation}\label{transport:first-flow}
\dot y=-v_\sigma(p,y),\qquad y(0)=y_0.
\end{equation}
Since \(J_j\dot y=-e_\sigma\), differentiation of
\eqref{transport:coordinate-map} shows that the selected slacks and
\(\delta=r+w\) are constant, while
\(\sigma(t)=\sigma_0-t\) and \(\tau(t)=\tau_0\).
We stop at time \(\sigma_0\); if \(\sigma_0=0\), this first segment
is empty.  From its endpoint, solve
\begin{equation}\label{transport:second-flow}
\dot y=-v_\tau(p,y)
\end{equation}
for time \(\tau_0\).  On the second segment, \(\sigma=0\) and
\(\tau(t)=\tau_0-t\), with the same twenty slacks and \(\delta\)
constant.  Neither gap crosses zero before its prescribed stopping
time.

For completeness, these stopping times are reached by actual solutions.
The vector fields are smooth wherever \(J_j\) is invertible, and this
open set contains \(Q_j\).  On either segment, for total elapsed time
at most \(T=\sigma_0+\tau_0\), the displacement from \(\xi_0\)
is at most \(3T\) in infinity norm.  Therefore every part of a
solution before its prescribed endpoint lies in the compact inner box
\[
K=\{\xi:\iin{\xi-c_j}\le d_0+3T\}\subset\operatorname{int}Q_j.
\]
This estimate first holds up to any possible first exit and precludes
that exit by \eqref{transport:strict-start}.  The usual continuation
theorem for a smooth finite-dimensional vector field then rules out a
finite maximal existence time before either prescribed endpoint:
the trajectory remains in a compact subset of the field's domain.
The parameter coordinate remains exactly fixed during this argument.

The twenty selected physical slacks remain nonnegative because they
retain their initial values.  The same is true of \(\delta\), while
the two gaps remain nonnegative by their explicit evolution.  The
other 81 source slacks remain strictly positive by
\eqref{transport:margin-main}.  Thus the complete entry and pivot-stage
bounds hold simultaneously for \(M(\xi(t))\) at every point of both
segments.  Pivot ties are allowed in these bounds; the construction
does not exchange unequal pivots or divide by a vanishing pivot.

At the endpoint both gaps vanish, so the actual tail is
\(\left(\begin{smallmatrix}r&r\\r&w\end{smallmatrix}\right)\).
This is a complete \(X\) source, whether or not \(\delta=0\).
Finally, the minus signs in \eqref{transport:first-flow} and
\eqref{transport:second-flow}, together with
\eqref{transport:direction-main}, give \(\dot F\ge1/10\) on each
nonempty segment.  Integration gives the stated height gain, and
Proposition~\ref{prop:X} bounds the endpoint by \(\alpha\).
\end{proof}

For a proper negative-diagonal source the conclusion reads
\(F\le\alpha-(2r-s-t)/10\).  More generally, every physical source
in a closed inner box \(\iin{\xi-c_j}\le1/10000\), with nonnegative
gaps, satisfies \eqref{transport:strict-start}, since
\[
\frac1{10000}+3\left(\frac1{10000}+\frac1{10000}\right)
=\frac7{10000}<\frac1{1250}.
\]
The strict inequality in \eqref{transport:strict-start} is part of the
terminal rule.  A source box is accepted by this rule only when its
entire enclosure satisfies that strict distance-and-time bound.
Equality at this outer guard is not accepted by a limiting argument;
it remains with the other sources to be subdivided or covered by a
separate valid terminal.  Zero gaps themselves are permitted, as the
proof omits the corresponding zero-length segment.  These four local
regions are not asserted to cover the whole proper negative-diagonal
domain.  Further safe transformations obey the same source and
height-direction requirements.


\subsection{Complete composition}
\begin{lemma}[Mixed threshold and algebraic coverage]\label{lem:mixed}
Let \(\gam<\alpha\). Let a rational root \(Q_0\) contain a canonical
image of every hypothetical source with \(F>\alpha\), with
height at least that of the source.
Suppose a finite acyclic rooted proof graph covers \(Q_0\) by exact
subdivisions, qualified contractions, and source-preserving coordinate
changes that do not decrease the height of a hypothetical source.
Suppose every terminal establishes either:
\begin{enumerate}
\item no complete source with \(F\ge\gam\) exists in its region; or
\item every complete source in its region has \(F\le\alpha\).
\end{enumerate}
If every child is accounted for and each terminal source is bound
to its declared root path, then no original source has \(F>\alpha\).
\end{lemma}
\begin{proof}
A hypothetical source has a canonical image in \(Q_0\) with no
smaller height, still exceeding \(\alpha>\gam\).
At a subdivision it belongs to a closed child; overlaps on boundaries
are harmless. A qualified contraction has already certified its
discarded region impossible or safe. A coordinate change maps the
same source into its declared image at no smaller height.
Following these relations through the finite complete graph reaches
a terminal and contradicts either terminal assertion.
\end{proof}

Previously accepted subproofs may be composed by this lemma, but only
on the same mathematical sources. The input includes the mathematical domain, not only its dimension
or a previously recorded success label. The implementation binds model,
parent, path, endpoints, transformation prefix, rule implementation,
and proof payload. The final traversal reconstructs the original
root boxes and verifies every grafted subproof's source.
Hashes establish byte identity; mathematical validity comes from
the accepted subproof and its declared rule.

\section{The complete negative-diagonal certificate}\label{sec:Bcomplete}

\begin{proposition}[Certified complete negative-diagonal bound]
\label{prop:B}
Every complete real source with \eqref{eq:BD},
\(p,k,r>0\), and \(0<s,t<r\), satisfies
\[
F=r+st/r\le\alpha.
\]
\end{proposition}

The computation uses the full seventeen-variable root \eqref{eq:Broot}
with Theorem~\ref{thm:fibre}'s exact framework semantics.
Its model contains every hypothetical \(F>\alpha\) source after
the three actual paths; it imposes no selected tail-contact equation
on all canonical images. Supplementary Materials~S5--S7 specify the
model, acceptance rules, and complete certificate data.

\subsection{Base certificate and remaining source domains}
Exact rational verification of the base certificate accepted
\(749693\) nodes: \(374846\) splits, \(374839\) contradiction terminals,
four safe terminals, and four open terminals. Its \(750\) enclosed
continuation certificates had no internal open terminal.
The four remaining domains are closed dyadic boxes reconstructed
from the original seventeen-dimensional root. Denote them by
\(Q_1,Q_2,Q_3,Q_4\), in the order specified in Supplementary Material~S7.
The base certificate proves the required bound outside their union.

Complete certificates for \(Q_1,Q_3,Q_4\) cover, respectively,
\(111\), \(159\), and \(56\) constituent sources.
The certificate for \(Q_2\) has \(240\) constituent sources;
\(239\) are covered by its component certificates and the remaining
one by the local completion described below. Every constituent
source retains its complete ancestor transformation prefix.

\subsection{Local completion}
The remaining source in \(Q_2\) is covered by a binary subtree with
\(54\) distinct, prefix-free terminal paths. Their union exhausts
that subtree. The terminal justifications are
\begin{center}
\begin{tabular}{@{}lr@{}}
\toprule
Terminal justification&Count\\
\midrule
Contradiction certificates&48\\
Height and feasible-transport certificates&3\\
Canonical-source threshold exclusions&3\\
\midrule
Total&54\\
\bottomrule
\end{tabular}
\end{center}
Each canonical-source threshold exclusion covers both qualified
charts of its declared source. The three height or transport
certificates establish safety at \(\alpha\).
The mathematical verifiers and source-prefix checks were executed
for this local completion. It covers one constituent source of
\(Q_2\); full coverage of \(Q_2\) also uses the other \(239\) sources.

\subsection{Composition at the original root}
The local completion is attached to its exact source in \(Q_2\),
preserving the already verified complement. The four complete
certificates are then attached to the four open terminals of the
base certificate. The composition traverses the actual base-root
structure and reconstructs its entire set of open boxes.
For each attachment it checks the original model, root path,
box endpoints, and full transformation prefix against the
corresponding complete certificate.
The reconstructed boxes coincide with \(Q_1,Q_2,Q_3,Q_4\).
Thus every open base-root source is covered by a matching complete
proof, and the assembled certificate has no uncovered source.

The constituent-source counts above total
\(111+240+159+56=566\). These are internal obligations of the four
remaining parents, not a count of original parents. They discharge four
original-parent obligations; together with the previously covered
\(465\), they give the original-parent coverage count \(469/469\).
Appendix~\ref{certificate:details} makes this distinction explicit and
describes the source-by-source correspondence.
This count provides a consistency check; coverage follows from
the complete root traversal and exact source comparisons.
Supplementary Material~S7 records the paths, endpoints, certificate
identities, and composition results needed to check these bindings.

\begin{proof}[Proof of Proposition~\ref{prop:B}]
A hypothetical \(F>\alpha\) source has \(p>1\) by \eqref{eq:F4p}.
Its actual sign normalization and Theorem~\ref{thm:fibre}'s paths
produce a complete framework in \eqref{eq:Broot} with
\(\cB>\alpha\). The accepted base proof and the four full-parent
completions form the finite complete cover just described.
Each contradiction terminal excludes this high source, and each
safe terminal bounds its height by \(\alpha\). All original root
paths and transformed sources are accounted for.
Lemma~\ref{lem:mixed} therefore excludes the hypothetical source.
\end{proof}

\begin{remark}[Verification and composition]
The final composer reuses identified, previously accepted complete-root
and parent receipts. It verifies their producer and payload identities
and exact source bindings; it does not rerun every terminal's
mathematics. The base tree retains four open records, each discharged by an
attached complete certificate. The records retain their original
meaning in this composition.
Likewise a source safe at \(\alpha\) need not be infeasible at
the lower threshold \(\gam\).
\end{remark}

\section{Completion of the global proof}
\begin{proof}[Proof of Theorem~\ref{thm:main}]
The lower bound is \cite[Theorem~2.2]{ChenEdelmanUrschel2026};
Section~\ref{sec:alpha} describes the exact corresponding witness
in the archive.

For the upper bound, fix an arbitrary nonzero real matrix and
arbitrary legal complete-pivoting path. Normalize its maximum initial
entry to one and absorb the path's permutations into the matrix.
If a pivot vanishes, its whole trailing block vanishes.
Corollary~\ref{cor:firstfour} bounds all growth through the fourth
stage by \(4<\alpha\). It remains to exclude a fifth-pivot value
larger than \(\alpha\).

Proposition~\ref{prop:lift} gives the complete two-step lift.
Lemma~\ref{lem:saturation} gives a same-height complete source on
the \(D\) or \(X\) saturated boundary. Actual sign changes put every
\(X\) source into \eqref{eq:Xtail}, bounded by
Proposition~\ref{prop:X}.
For a \(D\) source, \eqref{eq:Dtail} with positive sign or a zero
arm has height at most \(r\le4\); a tied arm goes to \(X\).
Every remaining source belongs to \eqref{eq:properB}, so
Proposition~\ref{prop:B} applies.
No fifth pivot exceeds \(\alpha\).
Since the matrix and legal path were arbitrary,
\(\rho_5^\R\le\alpha\), proving equality.
\end{proof}

\section{Reproducibility and scope}\label{sec:reproducibility}
The proof can be read at two complementary levels. The analytic
reductions and local acceptance rules specify why the conclusion follows
from the stated finite checks. The published exact data make those
checks independently executable. Reading the proof and inspecting the
rules therefore exposes the mathematical obligations; it does not
replace checking every large certificate instance.

The computational proof separates discovery, verification, and
composition. Discovery constructs cuts, subdivisions, and feasible
transformations. Verification checks finite rational inequalities
and qualified analytic rules on complete sources. Composition
establishes that verified pieces cover the original domain with
matching source data. This separation allows component proofs to
be checked and reused independently.

The supplementary documentation is organized into eight modules,
listed in Appendix~\ref{app:artifacts}. It identifies the polynomial
and algebraic witness, analytic reductions, complete \(X\) certificates,
negative-diagonal model, terminal rules, and final coverage graph.
The accompanying source and dependency index specifies exact file
identities, source domains, verification entry points, and the scope
of each recorded run. Large certificate payloads are separate from
the readable documentation; full verification requires those payloads
and all referenced component dependencies.

The base-root verification used exact rational arithmetic and took
\(22815.35\) recorded seconds. The final local verification took
approximately \(184.37\) seconds. These are the times for their
respective verification scopes, not the total cost of discovery
or of checking the entire proof. Later parent and root compositions
reuse the identified component verification records while checking
exact source bindings. A fresh end-to-end verification therefore
requires execution of every necessary component verifier as well
as the final composition. Checking a file's hash alone establishes
identity, not mathematical validity.

A companion Lean development formalizes the algebraic constant and its
attainment, key analytic reductions, and a conditional final theorem.
The latter retains two explicit global safety hypotheses, for the
physical \(X\) domain and the original \(B\)-root capacity endpoints
(named \texttt{XGlobalSafety} and \texttt{RootEndpointSafety}, respectively).
The present proof addresses these obligations through the written
analytic arguments, exact certificates, and source-composition logic;
their complete discharge in Lean, including the remaining model and
coverage connections, is not claimed. The Python/C++ verifier
implementations are not formally verified. This is a computer-assisted
proof with partial Lean formalization. Component verifications and
final composition are distinguished in Supplementary Material~S8;
composition does not re-execute every earlier verifier.

\paragraph{Lean project cold rebuild.}
On 12 September 2026, the frozen first-phase Lean project was rebuilt
on Linux from an empty project-object directory, using Lean 4.30.0
and pinned third-party dependencies. The recorded rebuild passed
all 634 product modules, three additional example/audit modules,
and 47 named axiom checks. It used an explicit dependency-order
driver, with third-party caches reused and supplemented as needed.
The final theorem's two safety parameters were checked and remain
explicit. Supplement S8 records the source identity, successful
build route, and scope of the retained evidence.

\paragraph{Data and code availability.}
The companion Lean source, exact verifier code, certificate archives,
and step-by-step verification documentation are available at
\url{https://github.com/hkjtsgmc79-boop/rho5-proof}.
The fixed release \texttt{v1.0.1} provides the certificate downloads,
SHA-256 manifest, file sizes, computing requirements, and recorded
verification scopes:
\url{https://github.com/hkjtsgmc79-boop/rho5-proof/releases/tag/v1.0.1}.
The project cold rebuild is distinct from a fresh-machine rebuild of
Lean and all dependencies, and from a new end-to-end mathematical
replay of the large certificates. Zenodo archival will follow separately.

The theorem determines the maximum value and supplies an attaining
matrix. Classification of all maximizers, uniqueness modulo symmetry,
and the analogous complex-matrix value remain separate questions.

\section*{Acknowledgments}
AI models and agent systems played an important role in this work,
contributing substantially to the development and verification of the proof.

Qianli Ma led and coordinated the research program, managed its
progress and budget, and made the final decisions on the selection,
continuation, or termination of research approaches. He organized
the allocation of scientific tasks and responsibilities among the
research agents.

The models and tools used included GPT, Claude, Wujie AI AGENT,
and Deepseek Harness. The authors retain responsibility for the
mathematical claims, supporting evidence, and final manuscript.

\appendix
% Replacement appendix section; insert after the manuscript's \appendix.
\section{Proof modules and supplementary material}\label{app:artifacts}

Supplementary Materials S1--S8 organize the analytic and computational
dependencies by their mathematical role. The theorem statements in the
text specify the domains on which each component is used.

\begingroup
\small
\renewcommand{\arraystretch}{1.08}
\begin{longtable}{@{}>{\raggedright\arraybackslash}p{.21\linewidth}
                        >{\raggedright\arraybackslash}p{.75\linewidth}@{}}
\toprule
Supplementary unit&Mathematical content and role\\
\midrule
\endhead
S1: Constant and attainment&
The degree-\(61\) polynomial, exact isolating interval, identity with
the constant of \cite{ChenEdelmanUrschel2026}, and supplementary
algebraic reconstruction of an attaining matrix. The lower bound is
\cite[Theorem~2.2]{ChenEdelmanUrschel2026}.\\[1mm]

S2: Scalar bounds and complete coordinates&
The conditional three-pivot envelope, the first-four-pivot bound,
the complete two-step lift, and the same-height \(D/X\) saturation
proved in Sections~\ref{sec:scalar}--\ref{sec:saturation}.
The saturation statement used here requires positive fifth-pivot height.\\[1mm]

S3: Canonical representative and height&
Qualified representative existence, exact row-scale reconstruction,
and canonical localization and height certificates, established
independently of the final negative-diagonal bound.
Six complete exclusion trees contain \(228542\) nodes and \(114274\)
contradiction leaves. The subsequent \(152467\)-record localization
uses \(208274\) contraction certificates; its three terminal
boxes are covered by the algebraic height theorem.\\[1mm]

S4: Complete cross-saturated domain&
The four source domains and exact sign and transpose maps of
Proposition~\ref{prop:X}, including their shared boundaries.
The low-\(r\) component uses S3.
All statements concern complete real matrices; no rank, cofactor-sign,
or auxiliary entry assumption is added.\\[1mm]

S5: Negative-diagonal framework&
The three complete feasible paths, attained six-variable maximum,
and seventeen-variable root of Section~\ref{sec:fibre}.
Zero factors, capacity ties, and the proper negative-diagonal
range \(k\le2\) are retained.\\[1mm]

S6: Exact acceptance rules&
Rational interval and nonnegative-combination certificates, qualified
contractions, algebraic-height rules, and complete feasible transports.
Contradictions above \(\gam\) and safety at \(\alpha\) have the
distinct semantics of Lemma~\ref{lem:mixed}.\\[1mm]

S7: Complete negative-diagonal cover&
The accepted base proof, complete parent proofs, and exact source
bindings of Section~\ref{sec:Bcomplete}.
The final \(54\)-terminal source completes the last parent's
\(240/240\) responsibilities and the total \(469/469\) original
parents. Four source-bound completions cover the base tree's four
retained open records; no original source region remains uncovered.\\[1mm]

S8: Reproduction and dependency index&
Model and payload identities, acceptance records, executable
dependencies, and source-composition maps.
The final composition's \(357\) direct dependencies are supplemented
by their recursive dependencies; the direct list alone is not the
complete proof archive.\\
\bottomrule
\end{longtable}
\endgroup

The dependency direction is S3 to S4, and then S4 to the
cross-saturated safety rules used in S6--S7; the final
negative-diagonal cover is not an input to S3 or S4.
Counts identify accepted objects, while their rules and exact
source bindings establish validity and coverage.
An \(\alpha\)-safe region may contain feasible matrices.

Component acceptance and final composition are separate operations.
The latter reuses identified acceptance records and checks their
source bindings; it does not rerun every component's mathematics.
Accordingly, the completed composition is not claimed to be a new
from-scratch replay of the entire dependency chain.
The supplementary index distinguishes the recorded acceptance runs,
the checks performed during composition, and the materials required
for independent reproduction.


\section{Canonical representatives and the exact near-peak height}
\label{s3:details}
% Included section: s3_detail
% Insert beneath the enclosing appendix section labelled s3:details.
\subsection{The original source and the precise reduction inputs}
\label{s3:source-section}

This appendix separates three operations: choosing a representative of a
prescribed attainable height; reconstructing that representative from one
set of row-scale coordinates; and bounding the resulting height by a
constrained maximum argument. The first operation uses earlier reduction
results. The second and third are stated and proved below, with the finite
interval statements isolated as computational propositions. This ordering
prevents the final \(X\) or negative-diagonal bound from being used to
justify an earlier near-peak terminal.

All indices of the \(3\times3\) blocks and their vectors are
\(0,1,2\). The letter \(s\) in this appendix means \(p-1\), and
\(R\) means the scalar \(F/2-1\). They are not the tail coordinate
\(s\) of the negative-diagonal fibre or the real field \(\R\).
We write \(\Gamma=V-\beta Y\) for a cofactor, reserving the
threshold \(\gam\) of the main paper for \(4132517/10^6\).

\begin{definition}[Complete canonical source]\label{s3:canonical-domain}
Take real parameters
\[
 (t,E,\beta,V,T,X,Y,Z,p,k,r,w,a,b,g,h,v,j,q)
\]
and set
\begin{equation}\label{s3:physical-matrix}
\begin{gathered}
 u=(1,V,T)^T,\qquad x=(X,Y,Z)^T,\\
 \widehat v=(-g,h,v)^T,\qquad
 \widehat q=(-j,-1-\beta h,q)^T,\\
 D=\begin{pmatrix}
 k&-a&-b\\ tk&r-ta&r-tb\\ -Ek&r+Ea&w+Eb
 \end{pmatrix},\qquad
 S=D+x\widehat q^T,\quad O=S+u\widehat v^T,\\
 P=\beta\widehat v+\widehat q,\qquad L=px-u,\\
 M=\begin{pmatrix}
 1&-1&\widehat v^T\\
 \beta&p-\beta&P^T\\
 u&L&O
 \end{pmatrix},\qquad F=r-w.
\end{gathered}
\end{equation}
The source belongs to \(\mathcal K\) when
\begin{equation}\label{s3:outer-physical}
\begin{gathered}
 p,k,r>0,\quad |t|\le1,\quad0\le E\le1,\quad
 0<\beta,T,X,Y,Z\le1,\quad0<V<1,\\
 0\le a,b\le k,\quad |w|\le r,\\
 |p-\beta|\le1,\quad |L_i|\le1,\quad
 |\widehat v_j|\le1,\quad |P_j|\le1,\quad |O_{ij}|\le1,\\
 |\widehat q_j|\le p,\quad |S_{ij}|\le p,\quad |D_{ij}|\le k,
\end{gathered}
\end{equation}
for every indicated index, and all six cofactors
\begin{equation}\label{s3:six-cofactors}
\begin{gathered}
\delta=1-\beta X,\quad\Gamma=V-\beta Y,\quad\nu=\beta Z-T,\\
\tau=Y-XV,\quad\Delta=Z-TX,\quad\mu=VZ-TY
\end{gathered}
\end{equation}
are strictly positive. In particular \(P_1=-1\) is built into the
representation. The signs of \(g,h,v,j,q\) are not extra assumptions
in this definition.
\end{definition}

These inequalities retain the complete initial and successive Schur
bounds. Eliminating the first two pivots gives the receiver/core blocks
with pivots \(p,k\); eliminating \(k\) gives
\(\left(\begin{smallmatrix}r&r\\r&w\end{smallmatrix}\right)\).
The last signed pivot is \(w-r=-F\). In the positive high-value
regime no division at a zero pivot is made. Sources with an earlier
zero pivot have already been handled before entry to this chamber.

The following is the precise inherited representative reduction. Its
source is the target theorem of the canonical-wall reduction, rather
than an assertion that every source is already a wall point.

\begin{proposition}[Prescribed-target representative]
\label{s3:target-reduction}
Put \(\Lambda=1653/400\). Fix a seven-coordinate head
\(\eta=(t,\beta,V,T,X,Y,Z)\) for which a complete source in
\(\mathcal K\) can have height at least \(f\), where
\(\Lambda\le f<5\). There is a complete matrix at this same head,
with height exactly \(f\), satisfying
\begin{equation}\label{s3:target-contacts}
\begin{gathered}
 w=-r,\quad r=f/2,\quad pZ=1+T,\quad
 P_1=-1,\quad P_2=1,\\
 O_{10}=O_{11}=1,\quad
 (O_{02},O_{12},O_{22})=(-1,1,-1),\\
 S_{10}=p,\quad O_{20}=-1,\quad S_{00}<p,\quad P_0>-1.
\end{gathered}
\end{equation}
Its cofactor and physical inequalities remain those of
Definition~\ref{s3:canonical-domain}. It also satisfies the
high-value qualifications recorded in \eqref{s3:high-budgets} below.
Conversely, a target representative at this head is a complete
physical source of height \(f\).
\end{proposition}

\paragraph{Scope and proof ingredients of the inherited proposition.}
Fix \(\eta\) and \(\Lambda\le f<5\), and assume that the complete
fibre at this head contains a source with \(F\ge f\).
We first establish the required attainment in the full physical domain,
before imposing any strict endpoint qualifications. On its subset
\(F\ge f\), the inequalities \eqref{s3:outer-physical} give
\[
 r\ge F/2\ge f/2,\qquad k\ge D_{21}=r+Ea\ge r,
 \qquad p\le1+\beta\le2,
\]
and, using the stage and auxiliary-vector bounds at the same source,
\[
 k=D_{00}=S_{00}-X\widehat q_0
 \le p(1+X)\le2p\le4.
\]
Consequently \(p\ge f/4>1\) and \(k,r\ge f/2>2\).
Also \(0\le a,b\le4\), \(|w|\le4\),
\(|g|,|h|,|v|\le1\), \(|j|,|q|\le2\), and
\(0\le E\le1\). All inner variables are bounded.
The physical conditions are polynomial weak inequalities, apart from
the positive pivots, whose uniform lower bounds have just been proved.
The six strictly positive cofactors depend only on the fixed head.
Thus this nonempty high-value fibre is compact, and the continuous
function \(F=r-w\) attains its maximum \(m\ge f\).
It is also the maximum on the whole fibre, since omitted sources have
height below \(f\).

The wall assertion is a separate inherited input. V20, \S7,
gives a complete wall representative at the same head and height
\(m>33/8\), using V19's wall result on \(\Gamma\le\nu\)
and the C19 exclusion on the remaining head region.
V18, \S5, then gives a wall of height exactly \(f\).
For this latter step, V15, \S\S2 and~4, first normalizes
\(p=(1+T)/Z\) and raises \(k\) to the complete first-column
capacity, preserving the head, \(E\), and height. In the resulting
V16 height polygon, decreasing the wall coordinates
\((m/2,m/2)\) to \((f/2,f/2)\) preserves every inequality:
the two oblique inequalities have positive coefficients, and the
remaining height constraints are upper bounds. Physical realization
uses specifically the wall constructions of V16, \S4.2, together
with V15's middle-column formula (5.5) at the new \(r=f/2\).
It does not assume that every low-value polygon point is realizable.
All columns use the same head, \(E,p,k\).

It follows that the set of complete sources at this head satisfying
\(F=f\) and \(w=-r\) is nonempty. It is a closed subset of the
compact high-value fibre above, so \(E\) attains a minimum on it.
Denote this value by \(E_{\min}\). This optimization concerns the
prescribed target \(f\), which need not be a maximum. It takes place
in the full \(\mathcal K\), not in the open predicate
\eqref{s3:QUAL}. The positive-tail theorem of V13, \S3, gives
\(v>0\) since \(f>33/8\); at this target wall the exact identity
\[
 Eb-(r-2)=(1+O_{22})+(1-Z)+Z(1-P_2)+\nu v>0
\]
shows that \(E_{\min}>0\). Indeed \(b\le4\) gives
\(E>(f-4)/8\). V16, \S2, excludes high-value sources with
\(\ell:=\nu-E\delta\le0\), so \(\ell>0\).
The inherited bound \(Ek<1-T\), together with \(k>2\),
gives \(E<1/2\).
Hence the minimum is strictly inside the upper computational endpoint
\(\min\{1,\nu/\delta\}\). No strict cofactor or pivot condition
is lost in a limiting argument.

Normalize \(p,k\) at this minimizing \(E\), keeping height \(f\).
Write \(\mathcal R_{\rm mid}(E)\) for the complete middle-column
capacity, and \(\mathcal W(E)\) for the minimum of the three complete
wall tail capacities of V23. The joint wall test is
\(r\le\mathcal R_{\rm mid}(E),\mathcal W(E)\).
These capacities are continuous near \(E_{\min}\). The original
first-column capacities have positive affine denominators in \(E\);
their common minimum is continuous even at ties. The middle-column
denominators remain positive because \(t+E>0\),
\(Q_o=t\delta-\Gamma>0\), and
\(t\nu-E\Gamma=t\ell+EQ_o>0\).
The tail denominators \(\Delta+\tau\) and \(\ell+Q_o\)
are likewise positive. If \(\mathcal W(E_{\min})>r\), continuity
preserves the tail test after a small decrease in \(E\).
V17, \S3.1, proves that this decrease preserves high-value first-
and middle-column feasibility; adopting the new maximal first-column
capacity only enlarges the middle-column feasible range. The same
target could therefore be reconstructed at a smaller \(E\), a
contradiction. Thus \(\mathcal W(E_{\min})=r\).

V23, \S3, expresses the active tail-capacity slacks as positive
weighted sums of the slacks of this one complete matrix. It follows
that
\[
 (O_{02},O_{12},O_{22})=(-1,1,-1),\qquad P_2=1\ \text{or}\ b=k.
\]
The common first- and middle-column normalizations also give
\(O_{10}=O_{11}=1\) and
\[
\begin{gathered}
 S_{00}=p\ \text{or}\ S_{10}=p,\qquad
 P_0=-1\ \text{or}\ O_{20}=-1,\\
 O_{01}=-1\ \text{or}\ O_{21}=1.
\end{gathered}
\]
Together with the wall and head contacts, these specify the complete
representative set used in V24, \S2. Only after reaching this set do
we invoke its exceptional-contact exclusion at \(F=f\ge\Lambda\).
The six systems exclude \(S_{00}=p\) or \(P_0=-1\) there;
they make no such assertion about arbitrary nonrepresentative sources.
Their finite certificates are inherited inputs; the argument here
identifies the physical set to which they apply.

The remaining contacts are \(S_{10}=p\), \(O_{20}=-1\),
\(S_{00}<p\), and \(P_0>-1\). V24, \S6, writes the four
capacity gaps as positive weighted sums of these same matrix slacks,
yielding \(k=K_{24}<K_{13},K_{14},K_{23}\) in
\eqref{s3:four-capacities}. If \(P_2<1\), then \(b=k\),
and the only active tail capacity is
\[
 \frac{\mu+\Delta+\tau+\Omega K}{\Delta+\tau},
 \qquad \Omega=t\Delta+E\tau-\mu.
\]
V19, \S4.4, strictly excludes the other eight original first-column
capacities. Together with the three gaps just obtained, this makes the
complete \(K\) locally equal to \(K_{24}\). These strict gaps
persist by continuity. Since \(d=tX-Y>0\), V24, \S7, gives
\[
 \frac{d(\Omega K_{24})}{dE}
 =-\frac{N_{24}V\mu d}{[V(tZ+EY)]^2}<0,
 \qquad N_{24}=p\mu+Y(V+T).
\]
Decreasing \(E\) improves this active tail capacity, preserves the
other tail margins, and does not impair the middle column. Minimality
again gives a contradiction. Therefore \(P_2=1\).

The two vanishing tail payments now give \(E=L_D\) and the formula
for \(b\) in \eqref{s3:QUAL}. V18 proves \(a_T<2<r\),
so its denominator is positive. V24, \S\S7--8, establishes the
remaining endpoint qualifications and the five joint middle-column
tests. Its theorem T24-target also proves sufficiency: the tests
restore the full middle-column conditions, the strict comparisons
restore the common first-column capacity, and V19's twelve-capacity
recovery and the wall construction produce one complete matrix of
height \(f\). These analytic and algebraic arguments apply to real
heads and real targets; the executable checker's rational-input restriction
does not restrict the mathematical quantifiers. Every use of a
representative above preserves all seven head coordinates.

\paragraph{Domains of the inherited high-value budgets.}
The inequalities \eqref{s3:high-budgets} are consumed only in their
complete positive canonical domain. In particular, V19, \S2,
proves \(F+\beta<5\) and \(g<3-F/2\) for complete sources
with \(F>4\); the form \(g<3-r\) is used only on the wall.
V22's exclusion gives \(d>0\) for original sources with \(F>33/8\);
its \S6 gives \(tp+d-1>0\) with the normalized
\(p=(1+T)/Z\), and \(1+t+d<33/16\) by a common
zero-tail-receiver reconstruction and V13's exclusion.
These are necessary conditions at the fixed head, not conclusions
from the exceptional-contact tree or the later height bound.
They do not extend this argument to arbitrary signed matrices or to
an unrestricted cofactor-nonnegative enlargement.

For clarity, the exact endpoint qualifications used in this proposition
are included here. They can either be retained in these coordinates or
replaced by the equivalent scale predicate in
Proposition~\ref{s3:scale-equivalence}. For the fixed head put
\[
 p=(1+T)/Z,\quad s=p-1,\quad
 Q_o=t\delta-\Gamma,\quad \ell=\nu-E\delta.
\]
The head entrance includes
\begin{equation}\label{s3:head-entrance}
\begin{gathered}
 0<t\le1,\quad0<V<1,\quad0<T,X,Y,Z\le1,\\
 t>V,\quad\tau>0,\quad\Delta>\tau,\quad\mu>0,\\
 \max\{0,p-1,1/X-1,T/Z\}<\beta<
 \min\{1,V/Y,5-f\},\\
 d=tX-Y>0,\quad tp+d-1>0,\quad1+t+d<33/16.
\end{gathered}
\end{equation}
Thus \(Q_o=\delta(t-V)+\beta\tau>0\).
Define
\begin{equation}\label{s3:four-capacities}
\begin{aligned}
 K_{13}&=\frac{p\Gamma+X(V+\beta)}{\Gamma+\beta tX},&
 K_{14}&=\frac{p\mu+X(V+T)}{\mu+X(VE+Tt)},\\
 K_{23}&=\frac{p\Gamma+Y(V+\beta)}{tV},&
 K_{24}&=\frac{p\mu+Y(V+T)}{V(tZ+EY)}.
\end{aligned}
\end{equation}
With
\[
 a_T=1-Y+\Gamma(1+X)/\delta,\qquad
 c_T=1+Z+\nu(1+X)/\delta,
\]
the endpoint and its complete qualifications are
\begin{equation}\label{s3:QUAL}
\begin{gathered}
 E=L_D:=\frac{\nu(r-a_T)+Q_o(r-c_T)}{\delta(r-a_T)},\qquad
 L_C:=\frac{r-1-Z}{1+X},\\
 E>L_C,\quad E<1,\quad\ell>0,\quad
 k=K_{24}>2,\quad k<K_{13},K_{14},K_{23},\\
 b=\frac{\delta(r-a_T)}{Q_o},\qquad 1+X<b\le k.
\end{gathered}
\end{equation}
Here \(a_T<2<r\), so no denominator in \eqref{s3:QUAL}
vanishes. All four capacity denominators are positive on the declared
head. The three strict comparisons in \eqref{s3:QUAL} are essential.

The five remaining middle-column tests are also joint tests at this
same \((\eta,E,k)\). Set
\[
\begin{array}{ll}
 A_1=1+Y+t(1-X),& B_1=Q_o,\\
 C_1=1+Y+t(p-X),& D_1=\Gamma+t\beta X,\\
 A_2=\dfrac{t(1+Z)+E(1+Y)}{t+E},&
 B_2=\dfrac{t\nu-E\Gamma}{t+E},\\
 C_2=\dfrac{tk+E(1+Y)}{t+E},&
 D_2=\dfrac{E\Gamma}{t+E},\qquad H=s/\beta.
\end{array}
\]
Their denominators and \(B_1,B_2,D_1,D_2\) are strictly positive.
The tests are
\begin{equation}\label{s3:M5}
 \begin{gathered}
 A_2+B_2H\ge r,\qquad A_1+B_1s\ge r,\\
 \frac{B_2C_1+D_1A_2}{B_2+D_1},\quad
 \frac{B_1C_2+D_2A_1}{B_1+D_2},\quad
 \frac{B_2C_2+D_2A_2}{B_2+D_2}\ \ge r.
 \end{gathered}
\end{equation}
For every head with the declared entrance conditions, the target
reduction says that existence of a complete source of height at least
\(f\) is equivalent to \eqref{s3:QUAL} and \eqref{s3:M5} at
\(E=L_D\). A common realization is
\[
 g=s/V,\quad j=(tk-p)/Y,\quad
 h=\max\{0,(r-A_1)/B_1,(r-A_2)/B_2\},
\]
\[
 a=(r-1-Y+\Gamma h)/t,\quad
 v=(b-1-X)/\delta,\quad q=1-\beta v,\quad w=-r.
\]
The proof uses the exact complete-column reconstruction, not separate
optimal completions of different columns.

The particular inherited necessary bounds used subsequently, at a
high target representative, are
\begin{equation}\label{s3:high-budgets}
\begin{gathered}
 0<g,j,h\le1,\quad v,q,\ell>0,\quad
 1<p<1+\beta<2,\quad a<p,\quad k\le9/4,\\
 0<t\le1,\quad t>V,\quad\Delta>\tau,\quad
 X>1/(1+\beta),\quad Ek<1-T,\\
 F+\beta<5,\quad g<3-F/2,\quad
 d>0,\quad tp+d-1>0,\quad1+t+d<33/16.
\end{gathered}
\end{equation}
They apply to complete high-value representatives in this positive
chamber; they are not imposed as necessary conditions on arbitrary
signed \(X\) matrices.

\begin{definition}[Admissible target representative]
\label{s3:target-representative}
An admissible target representative is a source in
\(\mathcal K\), at a height \(\Lambda\le F<5\), with all
contacts \eqref{s3:target-contacts}. In using its high-value
necessary inequalities we retain the inherited qualifications
\eqref{s3:high-budgets}. No conditions \(t=X=Z=1\), \(b=k\),
or \(h=s\) are included in this definition.
\end{definition}

\subsection{A complete row-scale coordinate system}
\label{s3:scale-section}

For an admissible target representative with \(\Lambda\le F<5\), write \(F=2r=2(1+R)\), \(p=1+s\), and use
\((\beta,V,s,Y,Z,v,h,b,R)\) as coordinates. Define
\begin{equation}\label{s3:scale-aux}
\begin{gathered}
 T=(1+s)Z-1,\quad q=1-\beta v,\\
 \Gamma=V-\beta Y,\quad\nu=\beta Z-T,\quad\mu=VZ-TY,\\
 A=R+Y+\Gamma v,\qquad C=R-Z+\nu v,\\
 W=R(Y+Z)+\mu v,\qquad N=V(Y+Z)+s\mu,\qquad
 \lambda=\frac{N}{VW}.
\end{gathered}
\end{equation}
The complete reconstruction is
\begin{equation}\label{s3:scale-dictionary}
 \begin{gathered}
 t=A/b,\quad E=C/b,\quad X=(b-1-v)/q,\quad k=\lambda b,\\
 g=s/V,\quad j=(\lambda A-p)/Y,\quad
 a=\frac{b(R-Y+\Gamma h)}A,\quad w=-r.
 \end{gathered}
\end{equation}
Inserting these expressions into \eqref{s3:physical-matrix} recovers
one matrix. In particular the products appearing in different
columns retain the same \(\beta,p,k,u,x\).

To derive the dictionary, the tail contacts give
\(tb=A\), \(Eb=C\), and \(b=1+v+Xq\).
The first-column contacts give
\(Vg=s\), \(tk=p+Yj\), and \(Ek+Zj+Tg=1\).
Since \(YC+ZA=W\), elimination in these three equations gives
\(k/b=N/(VW)\). The equation \(O_{11}=1\) then gives the formula
for \(a\). These are exact identities with the same variables.

Here is the full domain needed for the converse reconstruction.
Define the fixed-coordinate predicate \(\mathrm{BASE}\) by
\begin{equation}\label{s3:BASE}
\begin{gathered}
 0<\beta<1,\quad0<V<1,\quad0<T,Y,Z\le1,\\
 1<p<1+\beta,\quad F+\beta<5,\quad0<v\le s,\\
 \Gamma,\nu,\mu>0,\quad0<g\le1,\quad\lambda\ge1,\\
 0<j<1-\beta g,\quad pY-V\le1,\\
 D_h:=A\nu-C\Gamma=\mu+R(\nu-\Gamma)>0,\\
 n_h:=A(R-Z)+C(R-Y),\qquad
 0<h_0:=n_h/D_h\le H:=s/\beta.
\end{gathered}
\end{equation}
Throughout this appendix's height range, \(R>1\). Under
\(\mathrm{BASE}\), \(q>0\), \(A>C>0\), and all denominators
in \eqref{s3:scale-dictionary} are positive once \(b\ge A\).
Put
\begin{equation}\label{s3:scale-interval}
\begin{gathered}
 B_0=2+(1-\beta)v,\qquad
 B_\tau=1+v+qY/V,\\
 \sigma_0=\lambda q-j
 =\frac{(V+T)(gR-v)}{W}>0,\qquad
 B_F=\frac{pq-j(1+v)}{\sigma_0},\\
 I=[A,\infty)\cap(1+v,\infty)\cap(-\infty,B_0]
       \cap(-\infty,B_\tau)\cap(-\infty,B_F).
\end{gathered}
\end{equation}
The five conditions respectively mean
\(t\le1\), \(X>0\), \(X\le1\), \(\tau>0\), and
\(S_{00}<p\). The interval may be empty or a valid singleton.
A singleton at a closed endpoint is retained; a point on either of
the last two open endpoints is excluded.

Define
\begin{equation}\label{s3:capacities-scale}
\begin{gathered}
 \mathfrak a=R(1+q)+Vv,\qquad\mathfrak c=V+\beta R,\\
 L(h)=\frac{rA+C(R-Y+\Gamma h)}{A\lambda},\\
 U_O(h)=\frac{A[B_0+(1+\beta)h]}{\mathfrak a+\mathfrak c h},
 \qquad
 U_S(h)=\frac{A[pq+(1+v)(1+\beta h)]}{\mathfrak a+\mathfrak c h}.
\end{gathered}
\end{equation}

\begin{proposition}[Complete scale fibre]\label{s3:scale-equivalence}
Under \(\mathrm{BASE}\), a complete representative with the
specified fixed coordinates and target height exists if and only if
\begin{equation}\label{s3:complete-fibre}
 h_0\le h\le H,\qquad b\in I,\qquad
 L(h)\le b\le\min\{U_O(h),U_S(h)\}
\end{equation}
has a solution \((b,h)\). Every such solution reconstructs a matrix
satisfying all the original physical bands and all six positive
cofactors. Its pivots are \(1,p,k,r,-2r\).
\end{proposition}

\begin{proof}
The defining identities give
\begin{align}
 A(1-O_{21})&=D_hh-n_h,\label{s3:middle-identities-a}\\
 Aq(1+O_{01})&=A[B_0+(1+\beta)h]
                  -b(\mathfrak a+\mathfrak c h),\label{s3:middle-identities-b}\\
 Aq(p+S_{01})&=A[pq+(1+v)(1+\beta h)]
                  -b(\mathfrak a+\mathfrak c h).\label{s3:middle-identities-c}
\end{align}
The inequality \(r+Ea\le k\) is \(L(h)\le b\).
The receiver inequality \(|-1-\beta h|\le p\) is \(h\le H\).
Together with \eqref{s3:scale-interval}, these prove necessity.

For sufficiency, \(b\ge A>C>0\) gives \(0<t\le1\) and
\(0<E<1\), and \(b\in I\) gives \(0<X\le1\),
\(\tau>0\), and \(S_{00}<p\). The other cofactors follow from
\(\Gamma,\nu,\mu>0\),
\(\delta=1-\beta X>0\), and
\(\Delta=Z-TX\ge Z-T>0\). The signs and the opposite sides of
the physical bands are not discarded. They follow as follows:
\begin{itemize}
\item The head bounds use \(p-\beta\in(0,1)\),
 \(L_2=1\), \(L_0=pX-1\in(-1,1)\), and
 \(L_1=pY-V\in(-1,1]\).
\item The first receiver column has \(0<g\le1\),
 \(0<j<1-\beta g<1<p\), and
 \(P_0=-\beta g-j\in(-1,0)\).
 Its exact contacts are \(S_{10}=p\), \(O_{10}=1\),
 \(O_{20}=-1\). Also \(S_{20}=Tg-1\in[-1,0]\).
 From \(k\ge r+Ea>r>2\) and \(Xj<1\),
 \(S_{00}=k-Xj>1\). Thus
 \(0<O_{00}=S_{00}-g<p-g<1\), using \(g=s/V>s\).
\item The middle receiver has \(0<h\le H<1\) and
 \(-p\le\widehat q_1<-1\), with \(P_1=-1\).
 Equations \eqref{s3:middle-identities-b}--\eqref{s3:middle-identities-c}
 give its two nontrivial lower bands. Their opposite inequalities
 hold from the reconstructed signs; in particular
 \(S_{01}=-a-X(1+\beta h)<0\) and
 \(O_{01}=S_{01}+h<1\).
 Furthermore \(S_{11}=1-Vh\in(0,1)\), \(O_{11}=1\),
 and \(S_{21}=r+Ea-Z(1+\beta h)>r-2>0\).
 Together with \(O_{21}=S_{21}+Th\le1\), this pays both sides
 of the remaining middle-column bands.
\item The tail has \(P_2=1\),
 \((O_{02},O_{12},O_{22})=(-1,1,-1)\), and
 \[ (S_{02},S_{12},S_{22})=(-1-v,1-Vv,-1-Tv). \]
 The condition \(v\le s=p-1\) pays all three \(p\)-bands.
\item The core bounds use \(a<p<2<k\), \(b\le k\),
 \(tk\le k\), \(Ek<k\), and \(r+Ea\le k\).
 Here \(a>0\) follows from \(R>1\ge Y\), and
 \(a<p\) from \(S_{01}\ge-p\).
 The remaining entries satisfy
 \(D_{11}=1+Y-\Gamma h\in(0,2)\),
 \(D_{12}=1-Yq-Vv\in(-1,1)\), and
 \(D_{22}=-r+C\in(-r,0)\).
\end{itemize}
All entries are those of the same reconstructed matrix. These
identities pay the complete bands, not only a subset of necessary
inequalities. They also recover the strict first-column capacity
comparisons in \eqref{s3:QUAL} from the first-column residual
identities and \(S_{00}<p\), \(P_0>-1\).
\end{proof}

\paragraph{Why no division or endpoint is hidden.}
The interval \(I\) in \eqref{s3:scale-interval}, rather than its
closure, is used in the equivalence. The quantities
\(V,Y,q,A,\lambda,D_h,\sigma_0,\mathfrak a+\mathfrak c h\)
are positive in their stated domains. A zero denominator outside
this domain is not an infeasibility proof: such a point has not met
the declared qualification. Later, on the certified near-peak box,
these qualifications have uniform margins, so its closure does not
introduce zero-cofactor or excluded-endpoint cases.

\subsection{Common necessary inequalities and the finite localization}
\label{s3:localization-section}

Put \(\epsilon=R-1>0\). Two slacks of the same physical completion
have numerators
\[
 \mathcal T=(s-Vv)\mu-\epsilon V(Y+Z)=VW(\lambda-1)\ge0,
\]
\[
 \mathcal P=\epsilon V[(Y+Z)-v(\nu-\Gamma)]-(s-Vv)D_h
           =VW(1-j-\beta g)>0.
\]
Their sum is
\begin{equation}\label{s3:gamma-nu}
 \mathcal T+\mathcal P=(sR-Vv)(\Gamma-\nu).
\end{equation}
Because \(v\le s\) and \(V<1<R\), the first factor on the
right is positive. Therefore \(\Gamma>\nu\) throughout the
qualified high-value completion domain. This is a common-slack
identity, not a sign inferred from the candidate matrix.

The two middle capacities satisfy
\[
 U_O-U_S=\frac{Aq(h-s)}{\mathfrak a+\mathfrak c h},\qquad
 U_O'>0,\quad U_S'<0,\quad L'>0.
\]
For example
\[
 U_O'=\frac{2Aq(R-V)}{(\mathfrak a+\mathfrak c h)^2},\qquad
 U_S'=-\frac{Aq[\beta R(p-1-v)+V(p+1+v)]}
                  {(\mathfrak a+\mathfrak c h)^2}.
\]
It follows that every feasible \((b,h)\) satisfies
\(A\le b\le U_O(s)=U_S(s)\), without setting \(h=s\).
Define
\begin{equation}\label{s3:h-brackets}
\begin{gathered}
 d_1=1-V-\beta\epsilon,\quad
 a_h=\epsilon(1+q)-(1-V)v,\\
 d_S=V+\beta(\epsilon-v),\quad
 s_h=pq+1+v-\mathfrak a,\\
 h_A=a_h/d_1,\qquad h_S=s_h/d_S.
\end{gathered}
\end{equation}
Feasibility forces \(d_1>0\): indeed
\[
 B_0+(1+\beta)h-\mathfrak a-\mathfrak c h
       =d_1(h+v)-2\epsilon q
\]
and \(b\ge A\), \(b\le U_O(h)\) make the left side
nonnegative. Also \(d_S>0\), since \(v\le V\), \(\beta<1\),
and \(\epsilon>0\). The two middle inequalities then imply
\(h_A\le h\le h_S\). Equivalently, with
\[
 \mathcal M=(1-V)(s+v)-\epsilon[2+\beta(s-v)],
\]
one has \(\mathcal M\ge0\),
\(h_A=s-\mathcal M/d_1\), and \(h_S=s+\mathcal M/d_S\).

For reference we specify the necessary system used in the finite
localization. This makes the connection from the physical domain to
its rational certificate explicit. In the following table all
expressions are evaluated at the same six coordinates
\((\beta,V,s,Y,Z,v)\) and the same fixed \(R\); they are not
independently selected bounds. In addition to the definitions above,
put
\[
 u_0=rA+C(R-Y),\quad u_1=C\Gamma,
 \quad B_s=B_0+(1+\beta)s,\quad D_s=\mathfrak a+\mathfrak c s,
\]
\[
\begin{split}
 c_0&=NA^2B_0-VW u_0\mathfrak a,\\
 c_1&=NA^2(1+\beta)-VW(u_0\mathfrak c+u_1\mathfrak a),\\
 c_2&=-VWu_1\mathfrak c<0.
\end{split}
\]
Thus \(c_0+c_1h+c_2h^2=VW\Pi(h)\), where
\[
 \Pi(h)=A^2\lambda[B_0+(1+\beta)h]
          -(u_0+u_1h)(\mathfrak a+\mathfrak c h).
\]
The unconditional necessary inequalities are \(g_i\ge0\),
\(0\le i\le34\), with the following exact polynomial dictionary.
\begin{longtable}{@{}r p{.87\textwidth}@{}}
\toprule
\(i\)&\(g_i\)\\
\midrule
\endhead
0--6&\(\beta-s,\ T,\ \nu,\ \Gamma,\ \mu,\ s-v,\ V-s\), in this order.\\
7&\(B_0-A\).\\
8&\((s-Vv)\mu-\epsilon V(Y+Z)\).\\
9&\(d_1\).\\
10&\((1-V)(s+v)-\epsilon[2+\beta(s-v)]\).\\
11&\(D_h\).\\
12&\(\beta s-s^2-\epsilon\beta\).\\
13&\(3-2R-\beta\).\\
14&\((1+\beta)Y-V\).\\
15&\(1-pY+V\).\\
16&\(sD_h-\beta n_h\).\\
17&\(s_hD_h-n_hd_S\).\\
18&\(NAB_0D_h-VW(u_0D_h+u_1n_h)\).\\
19&\(NAB_0d_1-VW(u_0d_1+u_1a_h)\).\\
20&\(\mathcal P\).\\
21&\(Nq-(V+T)(sR-Vv)\).\\
22&\(pqVW-NqA+[Nq-(V+T)(sR-Vv)](A-1-v)\).\\
23&\(qY+V(1+v)-VA\).\\
24&\(NA^2B_sD_h-VWD_s(u_0D_h+u_1n_h)\).\\
25&\(AB_s-(1+v)D_s\).\\
26&\(c_1^2-4c_0c_2\).\\
27&\(9VW-4NA\).\\
28&\((9A-4u_0)D_h-4u_1n_h\).\\
29&\((1-T)VW-CN\).\\
30&\((3-r)V-s\).\\
31&\(17q-16R-16Vv+16v(1+\beta)\).\\
32&\(B_0(R+Vv)-A(1+v)\).\\
33&\((R+Vv-q)B_0+A(pq-1-v)\).\\
34&\(\Gamma-\nu\).\\
\bottomrule
\end{longtable}
The table is a necessary system, not the definition of a complete
physical source. Its rows follow by clearing positive denominators
in the scale conditions, the common bounds on \(h\), and
\eqref{s3:high-budgets}. To indicate all distinct sources: rows
0--17 use \(\mathrm{BASE}\), \(A\le b\le B_0\), the tail
capacity, and \(h_0\le H,h_S\); rows 18--19 use
\(L(h_0),L(h_A)\le B_0\); rows 20--23 use the strict first-column
bands and \(\tau>0\). Rows 24--25 use the shared middle bound at
\(s\), row 26 the discriminant of the same concave quadratic,
rows 27--30 the original pivot and head budgets, and rows 31--33
the three original geometric entrance inequalities. For the latter,
the identities are
\[
 d=\frac{R+Vv-A(1+v)/b}{q},\quad
 1+t+d=1+\frac{R+Vv}{q}-\frac{Av(1+\beta)}{qb},
\]
\[
 tp+d-1=\frac{(R+Vv-q)b+A(pq-1-v)}{bq}.
\]
The bounds \(A\le b\le B_0\) give the displayed necessary rows
with their stated directions. Row 34 is \eqref{s3:gamma-nu}.

There is one conditional row, not an additional unconditional
assumption:
\begin{equation}\label{s3:conditional-guard}
 g_{36}:=c_1+2c_2s\ge0
 \quad\Longrightarrow\quad
 g_{35}:=c_0+c_1s+c_2s^2\ge0.
\end{equation}
If a feasible \(h\le s\) exists, the nonnegative derivative at
\(s\) makes the concave quadratic increasing up to \(s\), so
\(\Pi(s)\ge\Pi(h)\ge0\). If \(h\ge s\), then
\(L(s)\le L(h)\le U_S(h)\le U_S(s)=U_O(s)\), with the same
conclusion. Any certificate row using \(g_{35}\) must first
establish \(g_{36}\ge0\) on its entire current box.

Set
\begin{equation}\label{s3:target-interval}
 f_0=\frac{4132517}{10^6},\quad f_1=\frac{103313}{25000},
 \quad R_0=\frac{2132517}{2000000},\quad R_1=\frac{53313}{50000}.
\end{equation}
The root enclosure for height \(f_0\) is
\begin{equation}\label{s3:localization-root}
 \mathcal B_0=[1/4,7/8]\times[1/16,1]\times[1/16,7/8]
          \times[0,1]\times[1/2,1]\times[0,7/8].
\end{equation}
Indeed the common middle and core bounds give
\[
 \epsilon+Ea\le\frac{T(\beta-s)}\beta
                    \le\frac{s(\beta-s)}\beta\le\frac\beta4.
\]
Thus \(s>\epsilon>1/16\) and \(\beta>4\epsilon>1/4\).
Also \(s\le V<1\), \(v\le s<\beta<5-f_0<7/8\), and
\(Z>1/(1+\beta)>1/2\). Hence every relevant six-tuple lies in
\(\mathcal B_0\). Its zero endpoints are only outer relaxations.

\begin{proposition}[Finite localization and upper-target exclusion]
\label{s3:finite-localization}
At \(R=R_0\), every point of \(\mathcal B_0\) satisfying the
35 unconditional inequalities and the conditional implication
\eqref{s3:conditional-guard} lies in
\begin{equation}\label{s3:Q}
\begin{split}
 Q={}&[0.617,0.6181]\times[0.7788,0.7795]
      \times[0.4525,0.4540]\\
   &\times[0.97352,0.97361]\times[0.99998,1]
      \times[0.1944,0.1950].
\end{split}
\end{equation}
All finite decimals here and below denote exact rational numbers.
At \(R=R_1\), the same necessary system is infeasible in the box
obtained from \(Q\) by replacing its last interval by
\([0.1944,0.19501]\).
\end{proposition}

\paragraph{Finite proof objects and acceptance rule.}
The first statement is the repaired complete localization certificate:
152,467 records, including 35,069 binary splits, 82,328 contraction
nodes, 3,338 direct contradiction terminals, 31,729 nonnegative
combination terminals, and three terminals wholly contained in \(Q\).
The contraction nodes contain 208,274 individually checked thin-slice
exclusions. The second statement is the repaired upper-target tree,
with 169 nodes and 85 contradiction terminals. Neither statement
uses the earlier, superseded localization or upper-target tree counts.

Their small mathematical acceptance rule is as follows. On a current
box use exact outward rational derivative and Hessian enclosures to
obtain affine upper bounds
\(g_j(x)\le b_j+\sum_i a_{ji}z_i\), \(|z_i|\le1\).
If \(g_j\ge0\), the resulting affine row is necessary. Nonnegative
rational weights \(w_j\) produce a contradiction when
\[
 \sum_jw_jb_j+\sum_i\left|\sum_jw_ja_{ji}\right|<0.
\]
The same inequality applied to an edge slice justifies a contraction.
Both children of each binary split cover their parent; the contracted
child contains every previously possible point. At a localization
terminal, containment of the entire current box in \(Q\) is
checked coordinate by coordinate. This proves the statement for
continuous real boxes; the rational mesh records enclosure endpoints,
not a discrete sample of possible source points. The complete finite
payloads and entry points are identified in \S\ref{s3:provenance}.

\subsection{Transport of the target and a compact nine-variable model}
\label{s3:transport-section}

The fixed-target statement alone would not suffice. Fix the old head
\(\eta\) of a high representative. Its tail identities give
\[
 tb=R+Y+\Gamma v,\qquad b=1+X+(1-\beta X)v,
\]
so
\begin{equation}\label{s3:target-transport}
 v(R)=\frac{R+Y-t(1+X)}{Q_o},\qquad
 Q_o=t(1-\beta X)-\Gamma>0.
\end{equation}
Choose its height-\(f_0\) representative using
Proposition~\ref{s3:target-reduction}. Its six-tuple lies in \(Q\).
Since \(b\ge A\), \(b\le B_0\), and \(X\le1\),
\[
 Q_o\ge (A/B_0)(1-\beta)-\Gamma>1/5.
\]
The last strict inequality is a rational enclosure on the entire
height-\(f_0\) box, not just its center. It follows that
\begin{equation}\label{s3:transport-size}
 0\le v(R)-v(R_0)
 \le5(R_1-R_0)=\frac3{400000}<10^{-5}
 \qquad (R_0\le R\le R_1).
\end{equation}
The coordinates \(\beta,V,s,Y,Z\) do not change with the target,
since \(s=(1+T)/Z-1\) is fixed by the old head.
If a source had height at least \(f_1\), it would have a target
representative at \(f_1\) with this same head. Equations
\eqref{s3:target-transport}--\eqref{s3:transport-size} would place it
in the upper-target exclusion box, a contradiction. This proves
\(F<f_1\) for every high representative, including those initially
specified above \(f_1\).
The comparison is between representatives at two targets; it does
not assert existence of a continuous physical flow in \(E\).

The bounds \(A\le b\le B_0\) and \(h_A\le h\le h_S\)
now put every height-\(F\ge f_0\) completion in
\begin{equation}\label{s3:nine-box}
\begin{split}
 \mathcal B={}&Q_{\beta,V,s,Y,Z}
 \times[0.1944,0.19501]_v\times[0.4513,0.4566]_h\\
 &\times[2.0741,2.0748]_b\times[R_0,R_1]_R.
\end{split}
\end{equation}
The relevant verified margins and enclosure statements are gathered
in the next proposition so that the inverse-function argument has
an explicit physical domain.

\begin{proposition}[Uniform qualifications]\label{s3:qualifications}
Throughout \(\mathcal B\), the rational expressions in
\eqref{s3:scale-aux}--\eqref{s3:h-brackets} have positive denominators,
and the following strict inequalities hold:
\begin{gather*}
 q,\Gamma,\nu,\mu,A,C,W,N,D_h,d_1,d_S,j,g,T>0,\\
 T<1,\quad g<1,\quad E<1,\quad0<v<s<\beta<1,\quad h<H,\\
 1-j-\beta g>1/1000,\qquad p-S_{00}>1/200,\\
 pY-V<1,\quad F+\beta<5,\quad X>0,\quad b>1+v,\\
 0<\beta,V,Y<1.
\end{gather*}
For any point in this box with \(X\le1\), also
\[
 \delta\ge1-\beta>0,\quad
 \tau\ge Y-V>0,\quad
 \Delta\ge Z-T>0.
\]
In the enlarged six-coordinate box and height interval, the same
rational enclosures give
\[
 A>2.0741,\quad B_0<2.0748,\quad
 h_A>0.4513,\quad h_S<0.4566.
\]
At the low target the transport denominator has the lower bound
\((A/B_0)(1-\beta)-\Gamma>1/5\).
\end{proposition}

These are 42 finite rational interval checks, including the matching
of the nine-dimensional Jacobian box with \eqref{s3:nine-box}.
The expression for \(1-j-\beta g\) is enclosed using
\(\mathcal P/(VW)\); this preserves the common-slack cancellation
in \eqref{s3:gamma-nu}. The three last cofactor lower bounds are
conditional on \(X\le1\), which will be supplied by an explicit
nonnegative slack, rather than being claimed on the whole outer box.

\subsection{Eight physical slacks and a uniform Jacobian}
\label{s3:slack-section}

On the domain of the reconstruction define the following eight
slacks, all with the feasible direction \(G_i\ge0\):
\begin{equation}\label{s3:eight-slacks}
\begin{aligned}
 G_1&=b-A,&G_2&=B_0-b,&G_3&=\lambda-1,\\
 G_4&=\lambda b-r-\frac{C(R-Y+\Gamma h)}A,\\
 G_5&=A[B_0+(1+\beta)h]-b(\mathfrak a+\mathfrak c h),\\
 G_6&=A[pq+(1+v)(1+\beta h)]-b(\mathfrak a+\mathfrak c h),\\
 G_7&=D_hh-n_h,&G_8&=1-Z.
\end{aligned}
\end{equation}
Their physical meanings, in the same order, are
\[
 \begin{gathered}
 t\le1,\quad X\le1,\quad b\le k,\quad r+Ea\le k,\\
 O_{01}\ge-1,\quad S_{01}\ge-p,\quad O_{21}\le1,\quad Z\le1.
 \end{gathered}
\]
For example \(G_1=b(1-t)\), \(G_2=q(1-X)\), and
\(G_3=(k-b)/b\); the identities
\eqref{s3:middle-identities-a}--\eqref{s3:middle-identities-c}
give \(G_5=Aq(1+O_{01})\), \(G_6=Aq(p+S_{01})\),
and \(G_7=A(1-O_{21})\). All these factors are positive.

\begin{lemma}[Exact feasible set in the localized box]
\label{s3:compact-model}
Let \(\mathcal C\) be the set of all complete scale representatives
at heights \(F\ge f_0\), with the source and qualifications of
Propositions~\ref{s3:target-reduction} and~\ref{s3:scale-equivalence}.
Then
\begin{equation}\label{s3:compact-set}
 \mathcal C=\{z\in\mathcal B:G_i(z)\ge0\ (1\le i\le8)\}.
\end{equation}
In particular, \(\mathcal C\) is compact and every member has
\(F<f_1\).
\end{lemma}

\begin{proof}
The forward inclusion is the localization and target-transport
argument above, followed by the physical interpretation of the slacks.
For the reverse inclusion, all the strict constant qualifications
in \(\mathrm{BASE}\) follow from
Proposition~\ref{s3:qualifications}; its non-strict tail capacity
\(\lambda\ge1\) is \(G_3\ge0\). To check explicitly the only
remaining condition involving \(h_0\), observe that \(R>1\),
\(Y<1\), and \(Z\le1\) imply \(n_h>0\). Thus
\(G_7\ge0\) gives \(0<h_0\le h<H\), using \(D_h>0\).
The constraints \(G_1,G_2\ge0\) give \(A\le b\le B_0\).
The other endpoints of \(I\) follow from \(X>0\),
\(\tau>0\), and \(p-S_{00}>1/200\); the cofactor bound for
\(\tau\) here uses precisely \(G_2\ge0\). Hence \(b\in I\),
with both strict endpoints respected. Finally \(G_4,G_5,G_6\ge0\)
are exactly \(L(h)\le b\le\min\{U_O(h),U_S(h)\}\).
Proposition~\ref{s3:scale-equivalence} now supplies every original
physical band of one complete matrix. The inherited high-value
necessary bounds apply to this complete representative as specified
in \S\ref{s3:provenance}.

All eight functions are continuous on a neighborhood of the closed
box because their denominators are bounded away from zero there.
The right side of \eqref{s3:compact-set} is therefore a closed
subset of a compact box. No previously strict cofactor or first-column
condition has been replaced by a boundary case. The independent
upper-target exclusion already proves \(F<f_1\) on this set.
\end{proof}

Write \(\zeta=(\beta,V,Y,Z,v,h,b,R)\), so that
\(z=(\beta,V,s,Y,Z,v,h,b,R)\) and
\(J=\partial(G_1,\ldots,G_8)/\partial\zeta\) at fixed \(s\).

\begin{proposition}[Uniform slack coordinates]\label{s3:jacobian}
For every \(z\in\mathcal B\), \(J(z)\) is invertible. In the
local coordinates \((s,G_1,\ldots,G_8)\),
\begin{equation}\label{s3:height-responses}
 \frac{\partial R}{\partial G_i}<-\frac1{250}
 \qquad(1\le i\le8).
\end{equation}
\end{proposition}

\paragraph{Finite rational verification of the proposition.}
The repaired certificate specifies a fixed rational matrix
\(C_0\), a rational vector \(y_0\), and positive rational radii
\(\rho_i\). With the rational midpoint of \(\mathcal B\) as
center, exact interval Hessians enclose every entry of \(J\)
by its central value plus its variation over all nine coordinates.
The verifier checks, for the entire box,
\begin{equation}\label{s3:jacobian-contraction}
 \mathcal E=I-C_0^T J^T,\qquad
 \|\mathcal E\|_\infty<\frac7{50}<1.
\end{equation}
Thus \(C_0^TJ^T\), and hence \(J\), are invertible.
Let \(e_R=(0,0,0,0,0,0,0,1)^T\). The same verifier bounds
the centered residual \(C_0^T(e_R-J^Ty_0)\) and proves
componentwise that
\[
 \left|C_0^T(e_R-J^Ty_0)\right|_i
       +\sum_j|\mathcal E_{ij}|\rho_j<\rho_i.
\]
Consequently the affine map
\(y\mapsto y+C_0^T(e_R-J^Ty)\) maps
\(y_0+[-\rho,\rho]\) strictly into itself, uniformly in \(z\).
Its unique fixed point solves \(J^Ty=e_R\), so
\(y_i=\partial R/\partial G_i\). In the stated order the exact
upper bounds \(y_{0,i}+\rho_i\) are
\[
 -\frac{31023639}{10^9},\quad
 -\frac{184717889}{10^9},\quad
 -\frac{47423419}{500000000},\quad
 -\frac{5888171}{25000000},
\]
\[
 -\frac{380657}{25000000},\quad
 -\frac{834537}{62500000},\quad
 -\frac{9580477}{200000000},\quad
 -\frac{4195499}{500000000}.
\]
Each is strictly smaller than \(-1/250\). These are interval
conclusions on the full nine-variable box, not gradients at its center
or a finite set of sampled points. The precise certificate data and
verifier are identified in \S\ref{s3:provenance}.

\begin{lemma}[A supercritical maximum has all eight contacts]
\label{s3:max-contacts}
If \(\mathcal C\) contains a point of height greater than
\(\alpha\), it has a maximizing point \(z_*\) with
\(f_0<F(z_*)<f_1\), and every such maximizing point satisfies
\(G_1(z_*)=\cdots=G_8(z_*)=0\).
\end{lemma}

\begin{proof}
Compactness gives a maximum. Under the stated hypothesis its height
is greater than \(\alpha>f_0\); the upper-target exclusion makes
it less than \(f_1\). This existence argument does not require an
additional numerical or rational witness.

Suppose \(G_i(z_*)>0\). The inverse-function theorem, using
Proposition~\ref{s3:jacobian}, supplies a local smooth inverse in
the eight slack coordinates at fixed \(s\). Decrease this one
slack by a positive amount \(u\), leaving the other seven fixed.
For sufficiently small \(u\), all eight slacks remain nonnegative.
At \(u=0\), \eqref{s3:height-responses} gives
\(dR/du>1/250\), so a short arc strictly increases height.
All the other strict qualifications persist by continuity. The
reconstruction theorem therefore makes this a feasible arc in the
original complete scale model, with height still above \(f_0\).

The artificial faces of \(\mathcal B\) are not additional
physical constraints. In constructing this short arc one works in
an open neighborhood on which the reconstruction denominators and
strict qualifications stay positive, and imposes only the eight
physical slacks. If a point of the arc appeared to leave
\(\mathcal B\), the already proved global localization, applied
to that original feasible point, would force it back into
\(\mathcal B\). Thus the arc belongs to \(\mathcal C\),
contradicting maximality. This proves all eight equalities.
\end{proof}

\subsection{The contact carrier and its exact elimination}
\label{s3:carrier-section}

The contact family is a consequence of the preceding maximum argument.
Indeed \(G_1=G_2=G_8=0\) gives \(t=X=Z=1\), whence
\(T=s\). The equations \(G_3=G_4=0\) give
\(b=k=r+Ea\). Also
\[
 G_5-G_6=Aq(h-s),
\]
so \(h=s\), and the last middle contacts are
\(O_{01}=-1\), \(S_{01}=-p\), and \(O_{21}=1\).
These contacts use the original shared matrix variables.

Put \(\epsilon=r-2=R-1\), and let
\(\mathcal L=(2-s)(1+s)\). We give the elimination from these
contacts to the carrier explicitly. The top middle contact gives
\(a=(1-\beta)s\); the bottom middle contact gives
\(\epsilon+Ea=s(\beta-s)\). Combining the latter with
\(b=r+Ea\) and the top tail contact gives
\[
 b=2+s(\beta-s),\qquad (1-\beta)v=s(\beta-s).
\]
The bottom tail contact is \(Eb=\epsilon+(\beta-s)v\).
Substitute these expressions into \((b-r)b=a(Eb)\).
After clearing only \(1-\beta>0\), the equation reduces to
\(2s(\beta-s)=\epsilon\mathcal L\). Thus
\begin{equation}\label{s3:carrier-dictionary}
\begin{gathered}
 \beta=s+\frac{\epsilon\mathcal L}{2s},\qquad
 v=\frac{s(\beta-s)}{1-\beta},\qquad
 b=2+s(\beta-s),\qquad a=(1-\beta)s,\\
 Y=1-\frac{\epsilon(s-v)}{s+v},\qquad
 V=1-\frac{\epsilon[2+\beta(s-v)]}{s+v}.
\end{gathered}
\end{equation}
For completeness, the last two formulas are the solution of the
middle and tail receiver equations
\[
 R-a=Y(1+\beta s)-Vs,\qquad b-R=Yq+Vv.
\]
Their determinant is \(v(1+\beta s)+sq=s+v>0\).
The remaining first-column relation, from
\(Eb+Zj+Tg=1\), \(j=(b-p)/Y\), and \(g=s/V\), is
\begin{equation}\label{s3:carrier-residual}
 \epsilon+(\beta-s)v+\frac{b-1-s}{Y}+\frac{s^2}{V}=1.
\end{equation}

To specify a compact integer polynomial, define
\[
 d_*=2s(1-s)-\epsilon\mathcal L,\qquad
 B_*=2s^2+\epsilon\mathcal L,
\]
\[
 Y_n=s(1-s)(1-\epsilon)+\epsilon^2\mathcal L,
 \qquad Y_d=s(1-s),
\]
\[
 V_n=2s^2(1-s)-2\epsilon d_*
       -\epsilon B_*s(1-s)+\epsilon^2B_*\mathcal L,
 \qquad V_d=2s^2(1-s).
\]
The subscript on \(d_*\) distinguishes it from the geometric
entrance quantity \(d=tX-Y\). Direct substitution in
\eqref{s3:carrier-dictionary} gives \(Y=Y_n/Y_d\) and
\(V=V_n/V_d\). Define
\begin{equation}\label{s3:carrier-poly}
\begin{split}
 G(s,\epsilon)={}&[2d_*(1-\epsilon)-\epsilon^2\mathcal L^2]Y_nV_n\\
 &-d_*(2-2s+\epsilon\mathcal L)Y_dV_n
       -4d_*s^4(1-s)Y_n.
\end{split}
\end{equation}
This polynomial has degree 10 in \(s\) and degree 7 in
\(\epsilon\). With \(\beta,v,b,Y,V\) given by the same
carrier dictionary, its exact relation to the residual is
\begin{equation}\label{s3:carrier-clearing}
 \epsilon+(\beta-s)v+\frac{b-1-s}{Y}+\frac{s^2}{V}-1
       =-\frac{G(s,\epsilon)}{2d_*Y_nV_n}.
\end{equation}
Here \(0<s<\beta<1\),
\(d_*=2s(1-\beta)>0\), \(Y_n=YY_d>0\), and
\(V_n=VV_d>0\). Thus no zero-denominator component has entered
the carrier equation: \eqref{s3:carrier-residual} is equivalent
to \(G(s,\epsilon)=0\) in the proved physical domain.

\begin{lemma}[Stationarity on the attained contact curve]
\label{s3:carrier-stationarity}
At the supercritical maximizing point of
Lemma~\ref{s3:max-contacts},
\[
 G(s,\epsilon)=0,\qquad G_s(s,\epsilon)=0.
\]
\end{lemma}

\begin{proof}
The same invertible Jacobian makes all eight coordinates \(\zeta\)
smooth functions of \(s\) when the eight slacks are held at zero.
For \(s\) in a two-sided open interval about its maximizing value,
all strict physical qualifications persist. The contact equalities
preserve the eight physical constraints, and the reconstruction
therefore gives a curve of complete physical representatives.
As in Lemma~\ref{s3:max-contacts}, global localization handles any
artificial face of the recorded box. Since the maximum is strictly
above \(f_0\), neither an auxiliary height endpoint nor an original
physical endpoint prevents this two-sided curve. Maximality gives
\(\epsilon'(s)=R'(s)=0\). Differentiating the identity
\(G(s,\epsilon(s))=0\) now gives \(G_s=0\).
This argument does not assume \(G_\epsilon\ne0\); the full
eight-slack Jacobian has already supplied the local parameter.
\end{proof}

\begin{proposition}[Exact elimination identity]\label{s3:resultant}
For the polynomial \eqref{s3:carrier-poly} and the normalization
of \(P_{61}\) in \eqref{eq:P61},
\begin{equation}\label{s3:resultant-identity}
\begin{split}
 \operatorname{Res}_s(G,G_s)
 ={}&\frac{\epsilon^{34}(\epsilon-2)(\epsilon-1)^2
             (\epsilon+1)^7(\epsilon^2-2\epsilon+2)^2}{131072}\\
 &\hspace{30mm}\cdot P_{61}(4+2\epsilon).
\end{split}
\end{equation}
\end{proposition}

This is a polynomial identity over \(\Q\), checked by computing
the resultant and verifying that the cleared polynomial difference
is identically zero. Its degree in \(\epsilon\) is 109; the
factor outside \(P_{61}\) has degree 48. The carrier dictionary
and \eqref{s3:carrier-clearing} are checked separately, so the
identity is connected to the complete physical maximum before
root identification is used.

At the maximum,
\[
 \frac{132517}{2000000}<\epsilon<\frac{3313}{50000}.
\]
Every factor outside \(P_{61}\) in
\eqref{s3:resultant-identity} is nonzero there. The two equations
of Lemma~\ref{s3:carrier-stationarity} imply that the resultant
vanishes, and therefore \(P_{61}(F)=0\).
Section~\ref{sec:alpha} has already proved that \(\alpha\) is
the unique root of this polynomial in \((4,5)\), which contains
\([f_0,f_1]\). Consequently the maximizing height must be
\(\alpha\), contradicting the assumed supercritical point.
The representative reduction then proves
Theorem~\ref{s3:canonical-height-main} for every source in
\(\mathcal K\). Neither a stationary-point calculation nor the
root-isolation certificate alone supplies this global bound.

\subsection{Source dependencies and the finite proof payloads}
\label{s3:provenance}

The reproducibility unit for the finite statements in this appendix is
\path{rho5_v28_exact.zip}, with internal root
\path{rho5_v28_exact/}. Its top-level report, Sections 1--5,
specifies the repaired localization and its interfaces. The
\path{RHO5_V26_REBUILT/REPORT.md} retained within it supplies
the mathematical argument in Theorem 6.1 and Sections 7--10.
The latter report's superseded tree counts and old Jacobian
contraction bound are not the data used in this appendix.

The responsibilities of the mathematical inputs are as follows.
\begin{enumerate}
\item The primary prescribed-target theorem is T24-target in
 Sections 7--8 of the theorem document in the original canonical
 representative archive identified in Supplement S3. Its nested
 sources give the wall-maximizer theorem (V20, Section 7),
 prescribed-height wall reconstruction (V18, Section 5; V16,
 Section 4.2), and common column reconstruction (V15, Sections
 2, 4--5). The compactness and minimum-attainment argument is
 given directly after Proposition~\ref{s3:target-reduction}.
 The wall choice and minimization in \(E\) precede the exceptional
 contact exclusions, which apply only in the complete representative
 set of V24, Section 2. Sections 6--8 of that source establish the
 common \(E=L_D\), \(k=K_{24}\) endpoint, its three strict capacity
 comparisons, and the five simultaneous middle-column tests.
 The V25 handoff in the localization package restates these inputs;
 it is not a substitute for their primary proofs.
\item The high-value inequalities \eqref{s3:high-budgets} are
 inherited only in their declared positive canonical representative
 domain and height range. In particular \(k\le9/4\),
 \(Ek<1-T\), \(g<3-r\), and \(F+\beta<5\) supply
 the common pivot and first-column budgets. The three V22 entrance
 inequalities \(d>0\), \(tp+d-1>0\), and
 \(1+t+d<33/16\) supply localization rows 31--33. These
 are not deductions from an arbitrary signed matrix or from the
 observed contact candidate. The first-column height payment is
 proved in V19, Section 2: it gives \(g<3-F/2\) and
 \(F+\beta<5\) before wall selection; \(g<3-r\) uses \(F=2r\)
 only after selection. V22, Sections 2--5, proves the head-only
 inequality \(d>0\) via a same-head wall maximum. Its Section 6.1
 proves \(tp+d-1>0\) after \(p=(1+T)/Z\) is imposed, and
 Section 6.2 proves the last head inequality by a common
 zero-receiver reconstruction and the safe-tail bound. These
 earlier V20/V22/V24 proof trees are separate dependencies;
 the localization payload itself does not re-execute them.
\item The V25 row-scale report in \path{sources/}, Sections
 2--4, supplies the exact shared reconstruction and its complete
 fibre equivalence. The dictionary, \(\mathrm{BASE}\), strict
 scale endpoints, and the necessity and sufficiency argument have
 been reproduced in \S\ref{s3:scale-section}. This appendix
 does not need the report's later finite elimination of \(\beta\).
\item The repaired V28 certificate supplies the continuous root-box
 coverage, the fixed-target transport margins, the upper-target
 exclusion, the nine-dimensional slack theorem, and the exact
 carrier/resultant connection. None of these statements uses the
 complete \(X\)-domain or negative-diagonal bound as an input.
 Root uniqueness and the definition of the exact constant are the
 separate result of Section~\ref{sec:alpha}.
\end{enumerate}

Supplement S3 identifies the complete replay entry point and the
separate localization, symbolic-interface, target-transport, Jacobian,
carrier, and resultant checks. Acceptance requires every localization
graft and all required modules. Discovery output or summary counts do
not replace that finite replay. The inherited representative reductions
remain separate mathematical dependencies: the canonical-height
certificate by itself does not prove the other sign chambers,
degenerate entrances, or the full matrix theorem.


\section{Data and coordinates for the two-gap transport}
\label{transport:details}
% Included section: transport_detail
% Body for the parent-supplied appendix section labelled transport:details.
\subsection{Complete coordinates and source inequalities}
\label{transport:source-detail}

The coordinates of \(\xi\) have the order in
\eqref{transport:coordinates-main}.  All entries of the vectors
\(u,x,v,q\) have indices \(0,1,2\).  Matrices \(D,S,O\) below use
these same indices.  The scalar coordinates \(A,B\) are entries of
\(D\), not the original input matrix.  Define
\begin{equation}\label{transport:actual-D}
D=\begin{pmatrix}
k&A&B\\
ck&r+cA&r-\sigma+cB\\
dk&r-\tau+dA&w+dB
\end{pmatrix},\qquad S=D+xq^T,\qquad O=S+uv^T.
\end{equation}
The associated real \(5\times5\) matrix is
\begin{equation}\label{transport:actual-M}
M(\xi)=
\begin{pmatrix}
1&-e&v^T\\
\beta&p-e\beta&(q+\beta v)^T\\
u&px-eu&O
\end{pmatrix}.
\end{equation}
Its successive leading Schur matrices, before the last scalar, are
\begin{equation}\label{transport:actual-stages}
\begin{pmatrix}p&q^T\\px&S\end{pmatrix},\qquad
D,\qquad
\begin{pmatrix}r&r-\sigma\\r-\tau&w\end{pmatrix}.
\end{equation}
These are algebraic identities for \(p,k\ne0\); elimination at the
fourth pivot additionally uses \(r\ne0\).  In particular, changing
the source coordinates in the transport changes one complete matrix
through \eqref{transport:actual-M}, with all its stage bounds retained.

For a scalar expression \(f\) with band \(|f|\le b\), write
\(f^+=b-f\) and \(f^-=b+f\).  The superscripts in the tables below
always denote these slacks, not the positive and negative parts of a
number.  The full polynomial list consists of the following bands,
with the identically zero slack \(D_{00}^+=k-k\) omitted:
\begin{equation}\label{transport:all-band-rows}
\begin{gathered}
e^\pm=1\mp e,\quad \beta^\pm=1\mp\beta,\quad
\mathrm{head}^\pm=1\mp(p-e\beta),\\
u_i^\pm=1\mp u_i,\quad x_i^\pm=1\mp x_i,\quad
v_i^\pm=1\mp v_i,\quad q_i^\pm=p\mp q_i,\\
L_i^\pm=1\mp(px_i-eu_i),\qquad
P_i^\pm=1\mp(q_i+\beta v_i),\\
D_{ij}^\pm=k\mp D_{ij},\qquad
S_{ij}^\pm=p\mp S_{ij},\qquad
O_{ij}^\pm=1\mp O_{ij}\qquad(0\le i,j\le2).
\end{gathered}
\end{equation}
There are 95 nonzero polynomials in
\eqref{transport:all-band-rows}.  The remaining nine are
\begin{equation}\label{transport:extra-rows}
r+w,\quad r-w,\quad p,\quad k,\quad r,\quad
\sigma,\quad\tau,\quad r-\sigma,\quad r-\tau.
\end{equation}
All 104 are required to be nonnegative for a physical source, with
\(p,k,r>0\).  On the boxes used here the strict positivity of these
three pivots follows from the stronger bounds below.  These conditions
give \(\mx{M}=1\), successive stage maxima \(p,k,r\), and legal
leading complete pivots \(1,p,k,r\), by
\eqref{transport:actual-stages}.  Since \(w<0\) on these boxes, the
last pivot is \(-F\) with \(F>0\).  No strict inequality between
competing maximal entries is required, so tied pivots are included.

For each \(j\), the components of \(\Phi_{j,p}\) are the twenty
slacks in Table~\ref{transport:slack-table}, in its displayed order,
followed by \(r+w,\sigma,\tau\).  The row index in that table is the
index within this selected list; it is not a numbering of all 104
source polynomials.  The Jacobian columns have the order
\begin{equation}\label{transport:jacobian-column-order}
\begin{split}
y={}&(k,r,w,A,B,c,d,e,\beta,u_0,u_1,u_2,x_0,x_1,x_2,\\
&\hspace{28mm}v_0,v_1,v_2,q_0,q_1,q_2,\sigma,\tau).
\end{split}
\end{equation}
Thus \(J_j\) is a \(23\times23\) matrix.  Its entries remain
polynomials in all 24 coordinates, including the parameter \(p\).

\begin{longtable}{rcccc}
\caption{The twenty selected source slacks in each local chart.}
\label{transport:slack-table}\\
\toprule
Index & \(j=0\) & \(j=1\) & \(j=2\) & \(j=3\)\\
\midrule
\endfirsthead
\toprule
Index & \(j=0\) & \(j=1\) & \(j=2\) & \(j=3\)\\
\midrule
\endhead
0 & \(e^+\) & \(e^+\) & \(e^+\) & \(e^+\)\\
1 & \(u_0^+\) & \(u_0^+\) & \(u_0^+\) & \(u_0^+\)\\
2 & \(x_0^+\) & \(x_0^+\) & \(x_0^+\) & \(x_0^+\)\\
3 & \(S_{01}^-\) & \(D_{01}^-\) & \(S_{01}^-\) & \(D_{01}^-\)\\
4 & \(O_{01}^-\) & \(O_{01}^-\) & \(O_{01}^-\) & \(O_{01}^-\)\\
5 & \(D_{02}^-\) & \(S_{02}^-\) & \(D_{02}^+\) & \(S_{02}^+\)\\
6 & \(O_{02}^-\) & \(O_{02}^-\) & \(O_{02}^+\) & \(O_{02}^+\)\\
7 & \(P_1^-\) & \(P_1^+\) & \(x_1^+\) & \(x_1^-\)\\
8 & \(D_{10}^+\) & \(D_{10}^+\) & \(L_1^+\) & \(L_1^-\)\\
9 & \(S_{10}^+\) & \(S_{10}^+\) & \(P_1^-\) & \(P_1^+\)\\
10 & \(O_{10}^+\) & \(O_{10}^+\) & \(O_{10}^-\) & \(O_{10}^+\)\\
11 & \(O_{11}^+\) & \(O_{11}^+\) & \(D_{11}^+\) & \(O_{11}^+\)\\
12 & \(O_{12}^+\) & \(O_{12}^+\) & \(O_{11}^+\) & \(D_{12}^+\)\\
13 & \(x_2^+\) & \(x_2^-\) & \(O_{12}^+\) & \(O_{12}^+\)\\
14 & \(L_2^+\) & \(L_2^-\) & \(P_2^-\) & \(P_2^+\)\\
15 & \(P_2^+\) & \(P_2^-\) & \(D_{20}^+\) & \(D_{20}^+\)\\
16 & \(O_{20}^-\) & \(O_{20}^+\) & \(S_{20}^+\) & \(S_{20}^+\)\\
17 & \(D_{21}^+\) & \(O_{21}^+\) & \(O_{20}^+\) & \(O_{20}^+\)\\
18 & \(O_{21}^+\) & \(D_{22}^-\) & \(O_{21}^+\) & \(O_{21}^+\)\\
19 & \(O_{22}^-\) & \(O_{22}^-\) & \(O_{22}^-\) & \(O_{22}^-\)\\
\bottomrule
\end{longtable}

\subsection{Exact rational centres}
\label{transport:centres-detail}
The following tables specify every coordinate of the four centres
exactly.  They are rational reference points for uniform estimates;
no optimality or exact contact identity is assumed at a centre.
The twelve common coordinates are
\begin{equation}\label{transport:common-centres}
\begin{gathered}
(c_j)_k=K_*:=\frac{518617093}{250000000},\qquad
(c_j)_r=R_*:=\frac{2066258539}{10^9},\\
(c_j)_w=-R_*,\\
(c_j)_p=\frac{726612331}{500000000},\qquad
(c_j)_\beta=\frac{617532517}{10^9},\\
(c_j)_{v_0}=-\frac{581690673}{10^9},\qquad
(c_j)_{q_0}=-\frac{159528331}{250000000},\\
(c_j)_e=(c_j)_{u_0}=(c_j)_{x_0}=1,\qquad
(c_j)_\sigma=(c_j)_\tau=0.
\end{gathered}
\end{equation}
For a compact exact table of the other twelve coordinates, define
\begin{equation}\label{transport:centre-table-constants}
\begin{aligned}
a_0&=\frac{10833981}{62500000},&
a_1&=\frac{47361589}{10^9},\\
a_2&=\frac{194787661}{250000000},&
a_3&=\frac{226612331}{500000000},\\
a_4&=\frac{194712659}{200000000},&
a_5&=\frac{194705107}{10^9},\\
a_6&=\frac{639940483}{500000000},&
a_7&=\frac{175952653}{200000000}.
\end{aligned}
\end{equation}
\begin{center}
\begin{minipage}{\linewidth}
\centering
\makeatletter\def\@captype{table}\makeatother
\caption{The remaining coordinates of the four rational centres.}
\label{transport:centre-table}
\begin{tabular}{ccccc}
\toprule
Coordinate & \(c_0\) & \(c_1\) & \(c_2\) & \(c_3\)\\
\midrule
\(A\) & \(-a_0\)&\(-K_*\)&\(-a_0\)&\(-K_*\)\\
\(B\) & \(-K_*\)&\(-a_0\)&\(K_*\)&\(a_0\)\\
\(c\) & \(1\)&\(1\)&\(-a_1\)&\(a_1\)\\
\(d\) & \(-a_1\)&\(a_1\)&\(1\)&\(1\)\\
\(u_1\) & \(a_2\)&\(a_2\)&\(a_3\)&\(-a_3\)\\
\(u_2\) & \(a_3\)&\(-a_3\)&\(a_2\)&\(a_2\)\\
\(x_1\) & \(a_4\)&\(a_4\)&\(1\)&\(-1\)\\
\(x_2\) & \(1\)&\(-1\)&\(a_4\)&\(a_4\)\\
\(v_1\) & \(a_3\)&\(a_5\)&\(a_3\)&\(a_5\)\\
\(v_2\) & \(a_5\)&\(a_3\)&\(-a_5\)&\(-a_3\)\\
\(q_1\) & \(-a_6\)&\(a_7\)&\(-a_6\)&\(a_7\)\\
\(q_2\) & \(a_7\)&\(-a_6\)&\(-a_7\)&\(a_6\)\\
\bottomrule
\end{tabular}
\end{minipage}
\end{center}
In particular, \((c_j)_r+(c_j)_w=0\), and
\(R_*-R=2065458539/10^9\).  A full outer box allows every coordinate,
including \(p\), to vary independently through its radius \(R\).
Only the evolution from an individual source fixes that source's
value of \(p\).

\subsection{Finite rational verification of the differential bounds}
\label{transport:rational-verification}
The companion proof-data table \texttt{two\_gap\_certificate.json}
contains the four \(23\times23\) rational matrices \(C_j\), as well
as the coordinate order, centres, and selected row labels.  Its data
schema is \texttt{V36\_TWO\_GAP\_FLOW\_V1}.  Rows of \(C_j\)
have the order of \(y\) in
\eqref{transport:jacobian-column-order}, and columns have the order of
\(\Phi_{j,p}\) in \eqref{transport:coordinate-map}.  The mathematical
role of this finite table is completely specified by the inequalities
below; the file name or its hash is not a substitute for those checks.

Each source slack has degree at most two.  Thus \(J_j\) is affine
in \(\xi\), as is \(E_j=I-C_jJ_j\).  For an affine entry
\(E_{j,ab}\), the exact whole-box bound used is
\begin{equation}\label{transport:affine-enclosure}
\sup_{\xi\in Q_j}|E_{j,ab}(\xi)|\le
|E_{j,ab}(c_j)|+
R\sum_{m=0}^{23}\left|\partial_{\xi_m}E_{j,ab}\right|.
\end{equation}
After summing across each row and taking the maximum, the right side
is \(\eta\) in \eqref{transport:contraction-main} for each of the
four charts.  The sum in \eqref{transport:affine-enclosure} includes
the \(p\) coordinate.  This is the explicit uniform parameter check,
rather than a computation at four parameter samples.

Since \(\iin{I-C_jJ_j}<1\), the Neumann series makes \(C_jJ_j\)
invertible.  Both matrices are square, so \(J_j\) is invertible and
\[
J_j^{-1}=(C_jJ_j)^{-1}C_j.
\]
For \(g\in\{\sigma,\tau\}\), put
\(v^0_{j,g}=C_je_g\) and \(n_{j,g}=\iin{v^0_{j,g}}\).
The exact inverse direction \(v_{j,g}=J_j^{-1}e_g\) obeys
\begin{equation}\label{transport:neumann-direction}
\iin{v_{j,g}-v^0_{j,g}}\le
\frac{\eta n_{j,g}}{1-\eta},\qquad
\iin{v_{j,g}}\le\kappa_{j,g}:=
\frac{n_{j,g}}{1-\eta}.
\end{equation}
The last two rows of \(J_jv_{j,g}=e_g\) give exactly
\[
((v_{j,\sigma})_\sigma,(v_{j,\sigma})_\tau)=(1,0),
\qquad
((v_{j,\tau})_\sigma,(v_{j,\tau})_\tau)=(0,1).
\]
Using \(F=r-w-\sigma-\tau+\sigma\tau/r\), differentiation gives
\begin{equation}\label{transport:height-derivative-exact}
D_yF\,v_{j,\sigma}
=(v_{j,\sigma})_r-(v_{j,\sigma})_w-1
 +\frac\tau r-\frac{\sigma\tau}{r^2}(v_{j,\sigma})_r.
\end{equation}
For \(v_{j,\tau}\), interchange \(\sigma\) and \(\tau\) in the
single-gap term.  On the full outer box,
\(r\ge r_{\min}>0\) and \(|\sigma|,|\tau|\le R\), even when a
point of the box is not physical.  Consequently the following rational
number is a valid upper bound in either case:
\begin{equation}\label{transport:height-upper-formula}
\begin{split}
U_{j,g}={}&(v^0_{j,g})_r-(v^0_{j,g})_w-1
 +\frac{2\eta n_{j,g}}{1-\eta}+\frac R{r_{\min}}\\
&\hspace{20mm}+\frac{R^2n_{j,g}}{(1-\eta)r_{\min}^2}.
\end{split}
\end{equation}

For clarity, all eight directional results can be specified by four
speed numerators and four height bounds.  Set
\[
D_v=9603274520441553881
\]
and
\begin{align*}
H_0&=-\frac{28753988041084852657186021690895856248996134979}
                 {102421779166357566006286717151929765002500000000},\\
H_1&=-\frac{12333488469648866145304175510636026684107958389}
                 {102421779166357566006286717151929765002500000000},\\
H_2&=-\frac{155271525823737135234142655316615127130663766999}
                 {204843558332715132012573434303859530005000000000},\\
H_3&=-\frac{4226128882854719189167870582514468491033177797}
                 {12802722395794695750785839643991220625312500000}.
\end{align*}
Substitution of the rational data into
\eqref{transport:neumann-direction} and
\eqref{transport:height-upper-formula} gives
Table~\ref{transport:direction-table}.  Every speed in the table is
less than \(3\), and each \(H_i\) is at most \(-1/10\).
\begin{table}[htbp]
\centering
\caption{Exact bounds for the eight inverse-coordinate directions.}
\label{transport:direction-table}
\begin{tabular}{cccc}
\toprule
Chart & Direction & \(D_v\kappa_{j,g}\) & \(U_{j,g}\)\\
\midrule
0 & \(\sigma\) & 14668693972000000000 & \(H_0\)\\
0 & \(\tau\)   & 24566153053000000000 & \(H_1\)\\
1 & \(\sigma\) & 12393729462000000000 & \(H_2\)\\
1 & \(\tau\)   & 10000000000000000000 & \(H_3\)\\
2 & \(\sigma\) & 10000000000000000000 & \(H_3\)\\
2 & \(\tau\)   & 12393729462000000000 & \(H_2\)\\
3 & \(\sigma\) & 24566153053000000000 & \(H_1\)\\
3 & \(\tau\)   & 14668693972000000000 & \(H_0\)\\
\bottomrule
\end{tabular}
\end{table}

It remains to verify the unselected physical inequalities on the
entire box.  If a source polynomial is
\(P(\xi)=\sum_{|\nu|\le2}a_\nu\xi^\nu\), its quadratic Taylor
expansion gives the rational lower bound
\begin{equation}\label{transport:quadratic-margin}
P(\xi)\ge P(c_j)-R\sum_{m=0}^{23}
 |\partial_{\xi_m}P(c_j)|
 -R^2\sum_{|\nu|=2}|a_\nu|\qquad(\xi\in Q_j).
\end{equation}
The 23 selected coordinate polynomials are distinct.  For every one
of the remaining \(104-23=81\) source polynomials, the right side of
\eqref{transport:quadratic-margin} is at least \(\mu\) in
\eqref{transport:margin-main}, in each chart.  This supplies the
physical interior margin used during continuation.  The selected
twenty band slacks and \(r+w\) do not need an interior margin:
their nonnegative initial values are conserved exactly, and the last
two selected coordinates are the explicitly decreasing nonnegative
gaps.  All comparisons in this subsection are comparisons of
rational numbers.

\subsection{Qualification of a closed source box}
\label{transport:terminal-detail}
For a closed rational enclosure
\(Q=\prod_{m=0}^{23}[L_m,U_m]\), define
\[
d_j(Q)=\max_m\max\{|L_m-(c_j)_m|,|U_m-(c_j)_m|\}.
\]
An immediately checkable sufficient guard for this transport is
\begin{equation}\label{transport:whole-box-guard}
L_r>0,\qquad L_\sigma,L_\tau\ge0,\qquad
d_j(Q)+3(U_\sigma+U_\tau)<R.
\end{equation}
Every complete physical source represented by such an enclosure
satisfies \eqref{transport:strict-start}.  The enclosure itself need
not consist of physical points; the conclusion is conditional on the
complete source inequalities, and does not infer simultaneous
feasibility from separate coordinate intervals.

For a proper negative-diagonal source described instead by
\((r,s,t)\), the conversion is the actual affine substitution
\(w=-r,\sigma=r-s,\tau=r-t\).  For example, the valid enclosing
interval for \(\sigma\) is \([L_r-U_s,U_r-L_s]\), and similarly
for \(\tau\).  This conversion may overestimate the image and may
therefore fail the sufficient guard; it does not justify replacing
the image by an unrelated smaller box.  If
\eqref{transport:whole-box-guard} is not satisfied, including its
equality boundary, this transport rule supplies no terminal there.
The complete proof graph must retain that region for another
qualified rule or for further subdivision.  Zero gap faces are
included whenever the strict guard holds, because a zero gap requires
no evolution in its corresponding direction.


\section{Source accounting and reproducibility}
\label{certificate:details}
% Included section: coverage_detail
% Body of Appendix D; the enclosing section supplies label certificate:details.
\subsection{Three levels of coverage}
\label{certificate:levels}

The numbers 54, 566 and 469 refer to different levels of the proof.
A terminal is a local region with an accepted conclusion. A constituent
source may require a tree of such terminals. An original parent
responsibility may in turn require many constituent sources.
In particular, the last 54-terminal anchor of
Section~\ref{certificate:worked} pays one constituent source, not 54
original parents.

The four residual original-root responsibilities have the following
bindings. The indices identify nodes in the retained root; the binary
paths identify those same nodes independently of the constituent
source numbering.
\begin{center}
\small
\begin{tabular}{@{}crlr@{}}
\toprule
Box & Original index & Original-root path & Constituents \\
\midrule
\(Q_1\)&338726&\texttt{10011000010101010}&111\\
\(Q_2\)&435003&\texttt{10111000001100011000010111011}&240\\
\(Q_3\)&563285&\texttt{1101000101111}&159\\
\(Q_4\)&675104&\texttt{1110101001111}&56\\
\bottomrule
\end{tabular}
\end{center}
Accordingly,
\[
\begin{gathered}
111+240+159+56=566\\
\text{constituent sources pay four original responsibilities},
\end{gathered}
\]
whereas the original-parent accounting is
\[
465\ \text{previously paid responsibilities}+4=469.
\]
There is no proposed bijection between the 566 constituents and the
469 original parents. For each of the four parents, the constituent
target-path set must equal its complete frozen source-path set, with
no omission or duplicate. Each target also retains its mathematical
acceptance and original binding evidence. Counts alone would not
establish any of these assertions.

The final anchor within \(Q_2\) contains 48 contradiction terminals,
three height or transport terminals, and three canonical-source
threshold exclusions. These are the same 54 positions under the
terminal semantics of Section~\ref{sec:certificates}. Combining them
proves safety at \(\alpha\), not that every region is empty at
\(\gam\). Local paths are interpreted relative to their frozen source
and its recorded transformations. In particular, a local path is not
silently concatenated with an original-root path to manufacture a
source binding.

\subsection{Why the original root still contains four open records}
\label{certificate:rootbytes}

The retained root has 749,693 nodes: 374,846 splits, 374,839
contradiction terminals, four \(\alpha\)-safe terminals and four
original open records. Its replay proves exactly this partial cover.
The four source-bound parent proofs above supply the missing
conclusions. The final overlay identifies each parent with the
corresponding original-root path and closed rational box, and checks
the identities of the accepted parent records. Applying the mixed
coverage lemma then gives an effective open count of zero.

The original root is preserved byte for byte. Its four open records
are neither deleted nor relabelled as contradiction terminals. This
separation permits an independent reader to replay the original tree,
replay the additional mathematical arguments and verify their precise
attachment points. A root replay that permits the four recorded open
positions is therefore an intermediate check, not acceptance of an
incomplete final proof.

\subsection{Locating the numerical example and its ancestors}
\label{certificate:exampledata}

The public release provides the 18.29 MB archive
\path{rho5-final-anchor-example.zip}. The input for
\eqref{certificate:actualmargin} is
\path{rho5_v53_exact/proof/TERMINAL_DETAILS.json}; the corresponding
\path{SOURCE_CERTIFICATE.json} and
\path{DEPTH19_COMPLETE_SOURCE_PROJECTION.json} bind it to the complete
depth-19 high source. The terminal-details file has SHA-256
\begin{quote}\small
\path{11d3e273fff5673cc9cb3ed9673860f388e685bd7ae3b769de0f89d39e20ecc1}.
\end{quote}
The archive contains all 49 weighted rows and all 67 coordinate
bounds used above, not only the two displayed rows.

The included restoration and joint-verification entry points replay
the complete 54-terminal anchor. Their successful result pays only
that anchor; the standalone whole-parent and whole-domain flags remain
false. The complete final-cover archive additionally contains the
frozen 239-source baseline, the completed parent record and the
original-root overlay. Supplement S7 identifies these objects and
their hashes. Thus the source-to-root chain for the example is
\[
\begin{gathered}
\text{depth-19 source and its contradiction}\\
\longrightarrow\ \text{complete 54-terminal anchor at }
\mathtt{10011101}\\
\longrightarrow\ \text{source 240 of 240 in parent 435003}\\
\longrightarrow\ Q_2\ \text{at its recorded original-root position}.
\end{gathered}
\]
The arrows denote complete source-cover compositions, not implications
from a terminal count or a receipt name.

\subsection{Independent verification and the supplementary map}
\label{certificate:verification}

The fixed public release is available at
\url{https://github.com/hkjtsgmc79-boop/rho5-proof/releases/tag/v1.0.1}.
Its index links all 24 original release assets; the archive inventory is
approximately 4.494 GB. The repository's
\path{docs/DOWNLOADS.md} gives sizes and archive identities,
\path{docs/VERIFY.md} gives the existing executable entry points, and
\path{docs/PROOF_MAP.md} distinguishes the written proof, exact
acceptors and the companion Lean statements. Large-file transport
parts reconstruct the original archive bytes and do not partition the
mathematical responsibility.

The supplementary index is organized by mathematical dependency:
\begin{center}
\begin{tabular}{@{}lp{0.79\linewidth}@{}}
\toprule
Part & Responsibility and verification material \\
\midrule
S1 & Exact constant, isolated root, attaining parameters and matrix.\\
S2 & Early pivot bounds, complete lift and saturation, including
      the zero-pivot qualifications.\\
S3 & Canonical representative, reconstruction and canonical height;
      the V24--V28 upstream dependencies.\\
S4 & All seven complete-\(X\) certificate components and their
      respective frozen entry points.\\
S5 & Six-variable elimination in the negative-diagonal domain and
      the source model used by the capacity certificates.\\
S6 & Necessary-row, height and transport acceptance rules, with
      their distinct thresholds and conclusions.\\
S7 & The preserved root, four parent bindings, the final anchor and
      source-cover composition.\\
S8 & The analytical, upstream, \(X\) and negative-diagonal archive
      collections, replay commands and scope of reproduction.\\
\bottomrule
\end{tabular}
\end{center}

There are three distinct checks. File hashes establish input identity.
Mathematical acceptors regenerate or verify necessary rows, rational
bounds, qualified analytic rules and local coverage. Source-cover
composition attaches these conclusions to the complete original
domain. The final composer consumes accepted records; rerunning it
alone does not rerun the mathematics behind those records. Its 357
direct file dependencies are also not a statement that the transitive
proof dependency closure has only 357 objects.

An independent complete reproduction therefore checks the upstream
analysis and \(X\) components, the original negative-diagonal root,
the four parent contributions, their height or transport mathematics,
and the final source bindings. Python assertion optimization must
remain disabled. Some historical configurations use absolute source
locations, so a new installation must supply an explicit path mapping
that preserves frozen identities and source checks. A tested,
location-independent command for the entire chain is not claimed.
The published component records remain evidence of their stated runs;
publication itself is not a new mathematical cold replay.

The exact Python and C++ acceptors are not asserted to be formally
verified programs. Their correctness, the written analytic arguments
and the source/domain composition remain part of this proof's trust
boundary. The companion first-phase Lean development records its
remaining global safety hypotheses explicitly; a successful build or
a certificate hash does not remove those hypotheses.


\clearpage
\bibliographystyle{plain}
\bibliography{references}
\end{document}
